
%\documentclass[oneside,final,14pt]{extreport}
\documentclass[12pt, a4paper, openany]{book}
\usepackage[left=2cm,right=2cm,top=1.5cm,bottom=1.5cm,bindingoffset=0cm]{geometry}
%\usepackage[koi8-r]{inputenc}
\usepackage[russianb]{babel}
\usepackage{vmargin}
\setpapersize{A4}
\usepackage[T2A]{fontenc}
\usepackage[utf8x]{inputenc} % more recent versions (at least>=2004-17-10)
%\usepackage[russian]{babel}
%\setmarginsrb{2cm}{1.5cm}{1cm}{1.5cm}{0pt}{0mm}{0pt}{13mm}
\usepackage{indentfirst}
\usepackage{nicefrac} % For comparison
%\usepackage{xfrac} % Works better with other fonts
%\usepackage[unicode, pdftex]{hyperref}
\usepackage{lettrine}
\usepackage[usenames]{color}
%\special{papersize=a5}
\usepackage{colortbl}
\usepackage[pagestyles]{titlesec}
\usepackage{xfrac} % Works better with other fonts
\usepackage[colorlinks=true,linkcolor=black,urlcolor=black,bookmarksopen=true]{hyperref}


% Настройка вертикальных и горизонтальных отступов
\titlespacing{\chapter}{0pt}{5pt}{5pt}
\titlespacing{\section}{\parindent}{4mm}{4mm}
\titlespacing{\subsection}{\parindent}{3mm}{3mm}

\newcommand{\anonsection}[1]{ \section*{#1} \addcontentsline{toc}{section}{\numberline {}#1}} 

\makeatletter %%%%% <---- Starting chapter without a pagebreak
\renewcommand\chapter{\par%
	\thispagestyle{plain}% \global\@topnum\z@
	\@afterindentfalse \secdef\@chapter\@schapter}
\makeatother %%%%% <---- Starting chapter without a pagebreak
\titleformat{\chapter}[display]
{\normalfont\bfseries}{}{0pt}{\Large}

\newpagestyle{mystyle}{
	\sethead[\thepage][][]{}{}{\thepage}	
}

\pagestyle{mystyle}

\sloppy
\begin{document}

	
	\begin{titlepage}
		
		\begin{center}
			\topskip0pt
			\vspace*{\fill}
			
			
			{\large\bf В.В. КСЕНОФОНТОВ\\}
			\ \\
			\ \\
			{\Huge\bf А ЕСЛИ В СЕМЬЕ ТЕЛЕВИЗОР И... ДЕТИ?\\}
			\ \\
			\ \\
			\ \\
			\vspace*{\fill}
			
			\vfill
			
			
			ИЗДАТЕЛЬСТВО <<ЗНАНИЕ>>
			
			Москва 1\ 9\ 9\ 3
			

		\end{center}
		
	\end{titlepage}
	
	\thispagestyle{empty} % выключаем отображение номера для этой страницы
	
	\newpage
	

\setcounter{secnumdepth}{0}
	
	\section[От автора]{\center \textbf{ОТ АВТОРА}}
	
Детище электроники — ТЕЛЕВИДЕНИЕ —  вошло в быт человека и неожиданно, и бурно. Особенно после Великой Отечественной войны. 

В 1950 году насчитывалось всего 4 тысячи владельцев  телевизоров. Их можно было бы назвать счастливчиками, если бы не одно обстоятельство, а именно — ежевечерне, помимо их желания, квартиры стихийно наполнялись соседями,  знакомыми, родственниками, которые регулярно заснимали места перед экраном. И это длилось по 3—4 часа. 

И вот — 1973 год. В стране работает около 130 телестанций, имеющих собственную программу. Телепередачами охвачена территория, на которой проживает более 170 миллионов  человек. Зрители знакомы с цветным телевидением, с  интервидением и, наконец, с космовидением. 

В зоне видимости телепередач трудно найти семью без  телевизора, так же трудно встретить пятилетнего ребенка, не  видевшего домашнего «кино» или не умеющего управлять этим  электронным аппаратом. 

Телевидение не только заметно повлияло на семейные  традиции и существенным образом изменило характер и структуру свободного времени людей, оно стало важнейшим источником получения знания. Одновременно ТВ (будем так его называть для краткости) поставило тысячи вопросов перед зрителями, психологами, врачами, философами, социологами,  литераторами, журналистами и педагогами. 

В 1963 году появилась замечательная книжка критика и публициста Владимира Саппака «Телевидение и мы». 

Это была первая попытка осмыслить ТВ как общественное явление, как искусство. 

С тех пор прошло десять лет. 

За этот промежуток времени в прессе появилось немало статей о ТВ, написаны диссертации. 

Но одной проблеме, связанной с телевидением, явно не повезло — проблеме влияния ТВ на воспитание детей. Хотя самый многочисленный и самый активный телезритель — это дети, подростки, юноши и девушки. 

Их многолетняя привязанность к ТВ невольно наталкивает на мысль о том, как повлияет телевидение на формирование подрастающего поколения? На поколение, которое с 3—4 лет начинает пополнять свой духовный багаж «телевизионной  продукцией» в значительных количествах и систематически.   Причем часто за счет заметного сокращения информации,  получаемой из других источников. 

Поскольку впервые предпринимается попытка поговорить о столь новом и сложном предмете, каким является проблема «Телевидение и дети», то мы были бы весьма признательны читателям, которые поделятся с нами своими мыслями и  опытом. 

Автору очень хотелось бы надеяться, что наши советы в соединении с вашим жизненным опытом, знаниями и здравым смыслом подскажут вам наиболее приемлемый вариант  использования телевидения в воспитании вашего ребенка. 

Итак, приглашаем вас на беседу, разговор, в котором мы хотим вместе с вами выяснить роль ТВ в жизни детей и  взрослых, а также в семейном общении; познакомить вас с разными подходами родителей к ТВ в воспитании детей; обратить  внимание на специфические особенности ТВ, его «силу» и  «слабости», на некоторые вопросы планирования передач и, наконец, дать ряд рекомендаций и советов. 

Проблема «ТВ и дети» — многосторонняя, поэтому автору приходилось не раз возвращаться к тому или иному вопросу, нарушая строгую последовательность изложения. 

Хотелось бы верить; однако, что это не помешает читателям извлечь из нашей работы определенную пользу. 
\newpage

	\section[Глава первая. Почему вы стали телезрителями? Какие невидимые связи установились между вашими детьми и телевидением?]{\center \textbf{ГЛАВА 1. ПОЧЕМУ ВЫ СТАЛИ ТЕЛЕЗРИТЕЛЯМИ? КАКИЕ НЕВИДИМЫЕ СВЯЗИ УСТАНОВИЛИСЬ МЕЖДУ ВАШИМИ ДЕТЬМИ И ТЕЛЕВИДЕНИЕМ? }}

Работа на телевидении воспитала у меня и  моих товарищей особое уважительное отношение к  телевизионным зрителям. 

Ведь приходится ежедневно обращаться с экрана к своим незримым друзьям, и они откликаются сотнями, тысячами писем. В них советы, пожелания, критика. 

Не было случая, чтобы зрители не поздравили нас с успешной передачей, даже если она была «неглавная». Мы постоянно чувствовали, как миллионы глаз телезрителей,  взыскательных и благодарных, следят за нашей работой и  способны заметить малейший промах и небрежность. Вот поэтому, предлагая новую передачу — плод раздумий, переживаний и труда целого коллектива, всегда испытываешь волнение. 

От имени авторов и режиссеров, от имени всех, кто создавал передачу, дикторы обращаются к зрителям «Дорогие друзья!» или «Дорогие товарищи!», «Дорогие советские воины», «Спортсмены», «Любители театра» и «Кинопутешествий...» и т. д. 

Эти обращения — не дань вежливости, а выражение  искренних чувств, желания быть полезными и понятыми теми, ради кого телевизионщики спускались в шахты, отправлялись в  труднодоступные места, репетировали новые спектакли,  разыскивали новые таланты. 

Вот и теперь по старой привычке мне хочется начать брошюру словами: 

Дорогие товарищи — родители, дедушки и бабушки! Давайте поговорим с вами о нашей  общей заботе — о ваших детях и нашем телевидении и в то же время о вашем телевидении и наших детях, ведь наши интересы полностью совпадают; и мы хотим и вы, чтобы наши советские дети выросли хорошими, счастливыми людьми, полноценными гражданами и строителями нового общества. 

А начнем наш разговор с небольшого воспоминания. С чего начиналось телевидение... в вашей семье? Очевидно, с  желания завести домашний кинотеатр, концертный зал, стадион. Раньше вы ограничивались патефоном и книгой, радио, ходили в клуб или кинотеатр, чтобы посмотреть новый фильм. Но ваши культурные потребности выросли. Вы хотите видеть и знать больше. Ну, а как быть, если вы живете не в Москве, а вам хочется, и даже необходимо, познакомиться со столичными спектаклями, выступлениями артистов эстрады? Вы хотите не просто знать результаты, а видеть ««наяву», в действии  советскую школу хоккея или фигурного катания. На наших глазах разворачивается космическая эпопея, люди штурмуют Луну, Марс, Венеру, и узнавать об этом только по фотографиям или документальным фильмам (а ведь кинофильм всегда «вчера», а «живое» телевидение — всегда «сейчас») — это уже, конечно, явно недостаточно. 

Тем более что «все вокруг давно имеют телевизоры». 

А кроме того, ваши дети время от времени просят вас: «давайте купим телевизор». А ваши родители-пенсионеры — разве они не хотят быть в курсе всех событий? Но здоровье ограничивает их возможности. И они тоже постепенно  убеждают вас купить телевизор. 

Купить телевизор — это значит в домашних условиях,  экономя время, силы и деньги, получать информацию,  развлечения, знания. Это расширение культурного кругозора и  одновременно отдых всей семьей. 

Разве не так? Именно эти аргументы приводят людей в магазин телевизоров. А иногда к ним добавляется еще один, — «чтобы сын поменьше болтался со своей дворовой компанией — пусть лучше набирается ума у телевизора». Резон? Для многих, пожалуй. И вы покупаете телевизор. И вот вы вдвоем-втроем, прихватывая иногда и детей, отправляетесь в магазин. Не спеша, приглядываетесь к телеприемникам с заманчивыми и разнообразными названиями, но, по существу, очень  похожими друг на друга. Часто не марка телевизора, и даже не стоимость его, а внешняя форма и цвет привлекают ваше внимание. Ну что же, и это важно и понятно. Наконец телевизор дома, установлена антенна, ожил экран. 

От мала до велика — вся ваша семья — несколько часов кряду смотрит передачи. И так — несколько дней, недель... Началась новая жизнь. Вы предстаете в новом качестве — вы уже телезрители. Это означает, что теперь по легкому  мановению вашей руки и легкого щелчка выключателя телевизора к вам в квартиру будут регулярно приходить люди и события, у вас появится много новых знакомых, появятся новые  привязанности, симпатии и... антипатии. Вы будете знать «все, что творится в мире», или, точнее, почти все, что волнует людей во всем мире. 

И ваше первое впечатление — будто к вам навсегда в дом пришел праздник и что множество вопросов, которые  возникали в семье до покупки «волшебного ящика»», наконец-то будут разрешены. 

Вы надеетесь, что теперь чаще будете всей семьей  находиться вместе, что это духовно объединит вас, укрепит  внутрисемейные отношения и т. п. 

Сознайтесь, действительно, подобное или нечто похожее пережили вы, купив телевизор? 

А что получилось через год-два, когда вы уже стали  телезрителем со стажем? Оправдались ли ваши надежды? Лучше всего об этом рассказывают потоки писем, которые идут на студии телевидения. В них благодарности, просьбы, советы,  рецензии. Приходят и жалобы на возникшие разногласия в связи с появлением в доме телевизора. Чаще всего они возникают из-за проблемы «Телевидение и дети». Дело порой доходит до самых крайних выводов: «Продадим телевизор!» Но давайте не будем спешить и попробуем внимательно разобраться в существе проблемы. 

Во-первых, как относятся к телевидению сами дети? 

Не появился ли у них, кроме вас и их сверстников, некто третий, кто занимает в их сердцах весьма заметное место? Да, появился. И этот «третий», разумеется, телевизор. Будем  терпеливы и послушаем, какие доверительные отношения  завязываются между вашими детьми и телевидением. 

«Дорогая тетя Нина! Мне уже шесть с половиной лет, я учусь в 1-м классе, хорошо читаю. 

Мне очень хочется, чтоб Вы или тетя Валя поздоровались со мной из телевизора. Зовут меня Костя Тютюнников,  пожалуйста». 

Задумайтесь, товарищи родители, почему так жизненно важно для Кости, чтобы диктор телевидения поздоровался с ним? Почему? Мальчик отыскивает бумагу и конверт, пишет «дорогой тете Нине» и просит... Что же он просит, чего ему не хватает? Игрушек? Книг? Очевидно, ему не хватает личного общения с доброй и спокойной тетей Ниной. 

Однажды, когда ребята узнали в одной из детских передач, что у известного циркового актера — Карандаша — скоро будет день рождения, студию телевидения завалили телеграммами, поздравительными открытками и маленькими самодельными сувенирами. 

Девятилетняя Марина Сомова писала: «Дорогой Карандаш, поздравляю с днем рождения. Вся наша семья, папа, мама, моя сестра Света и я, очень любим Вас. Желаем Вам счастья,  радости, не стареть и прожить еще сто лет. Горячо любящая Вас ученица 3-го класса средней школы № 28». 

Видимо, семья Марины, дружная, чуткая, телепередачи, в том числе и детские, смотрит коллективно, их обсуждают и вместе переживают. 

У Марины нет острой потребности в общении, как у Кости Тютюнникова, но у нее уже воспитана чуткость, отзывчивость к людям, благодарность за ту радость, которую они приносят  другим. 

Но продолжим наш разговор о причинах, побуждающих детей писать далекому и близкому телевидению. 

«Дорогая редакция! Я заметил, что в некоторых городах и деревнях не хватает асфальта. Я изобрел искусственный асфальт. Надо взять землю, черные чернила и канцелярский клей, а затем смешать и придать нужную форму, высушить. 

Я сейчас уже делаю цветной асфальт. 

А кому нужны искусственные чернила, надо взять воду, немного сажи и подлить черных чернил и все смешать, получаются очень хорошие чернила. Ваня». 

Мамы, папы, дедушки и бабушки! 

Сколько раз вы бранили своего маленького человечка за грязь, мусор, кляксы! А им двигала гражданская забота, чтобы и в деревнях были асфальтированные улицы. Сколько раз вы его не понимали! Вы думали только о том, чтобы он чернилами не испачкал скатерть, паркетный пол, стены. А ведь он не безделицей занимался. Он изобретал! И надо было бы ему помочь в его очень серьезном деле. 

Ну, а как же скатерть, стены? Если с уважением поговорить с ним, он, изобретатель, поймет вас. Да, поговорить, объяснить, а не ругать и не наказывать. Ни в коем случае не наказывать за изобретательство, не отбивать охоту к поиску, познанию,  исследованию, опыту. 

Часто детям кажется, что их на телевидении поймут скорее, чем дома, или даже в школе. Ведь там диктор никогда не  повышает голоса, не хмурится, не отмахивается, ссылаясь на занятость, а «если, что не так», то и тогда говорит  доброжелательно: «Дорогие ребята! Это делать не нужно», и спокойно объяснит почему. Видимо, поэтому, каких только сокровенных вопросов не встретишь в письмах на ТВ: 

«Дорогая редакция! Недавно мне поставили 2 по алгебре. Как мне выбраться из этого?» 

«А если по-настоящему, не в песне, то с чего же начинается Родина?». 

«Мы не знаем и никто нам не может ответить: почему мальчики и девочки в нашей школе всегда ссорятся, а в других всегда дружат?» 

Нину из Орехово-Зуева волнует совсем другое: «Если бы я была взрослой, я поехала бы помогать кубинскому народу строить новую жизнь или защищать Демократическую  Республику Вьетнам. 

Передайте привет, дорогое телевидение, Фиделю Кастро и вьетнамскому народу. Мы с ними. Мир и дружба». 

Вы не узнали, не почувствовали в этих письмах своих детей — любопытных, любознательных и беспокойных юных  граждан? 

Они хотят красиво танцевать, дружить, работать на  башенных кранах, помогать Кубе и Вьетнаму. 

Абсолютное большинство писем не вызвано какими-либо конкретными телевизионными передачами. Наоборот, они — трепетный отклик на события, происходящие рядом и в далеком мире. Это освоение действительности, умственное и  нравственное созревание, нелегкий путь становления личности. 

Они — наши дети — совсем не одиноки. Рядом с ними — всегда родители и близкие, сверстники и учителя. Однако  богатая и сложная жизнь с множеством загадок толкает ребят на поиски новых авторитетов в решении жизненно важных для них вопросов. И в подобной роли для одних выступает телестудия, для других — редакции радио и газет. 

Отражают ли письма отношение к телевидению тех ребят, которые не пишут и которых, как вы понимаете, неизмеримо больше? Да, отражают, это подтверждают социологические опросы. Обычно пишут или наиболее активные ребята, или те, кто надеется, что телевидение наилучшим образом ответит их «назревшим потребностям». Бывают письма, продиктованные, что называется, «криком души». Это те ребята, у кого нет близких друзей, кто по-настоящему одинок в семье или классе. 

Для подавляющего большинства детей телевидение стало и другом, и воспитателем, которого они выбрали добровольно. Когда одну девочку первоклассницу спросили, из кого состоит ее семья, она ответила, не без юмора, конечно: «В нашей семье мама, папа, бабушка и телевизор». 

Телевизор, особенно для малышей, — это что-то  одушевленное, близкое, домашнее, умное. 

Для ребят среднего возраста — нечто вроде друга, товарища в пути, он поможет, даст дельный совет и с ним нескучно. 

Для старших — советчик и источник знаний, путеводитель. 

Мы с вами рассмотрели только одну сторону проблемы  «Телевидение и дети» —это позицию детей. Она говорит о высокой степени доверия ребят к телевидению, о том, что у телевидения и юных зрителей установились добрые отношения, незримые, но прочные контакты. 

Как важно это понять! Как важно правильно оценить новые источники, влияющие на воспитание детей! 

Телевидение расширяет круг знаний, представлений о мире, телевидение — средство эстетического воспитания,  телевидение помогает в учебе, оно — средство развлечения и отдыха, организатор полезных общественных дел, телевидение  побуждает детей к коллективной и индивидуальной творческой деятельности, воспитывает нравственно и идейно. 

Тут все как будто бесспорно, все ясно — телевидению принадлежит огромная воспитывающая роль. Оно по своим целям и задачам — союзник семьи, школы и общественности. 

Однако мы нередко слышим вздохи родителей, бабушек и дедушек: «Да, телевизор — большое дело, но...» и дальше речь заходит о негативных сторонах ТВ. Журналисты и врачи,  педагоги и социологи тоже нет-нет да и обрушиваются на ТВ.  Вспомним, как писал Вл. Саппак: «Иллюзорны знания, не добытые трудом, пытливым поиском мысли, а воспринятые на слух,  наговоренные торопливыми лекторами с показом картинок...» «Жизнь скользит. Несется на тебя — и проносится мимо». «Жизнь мелькает, она не включила тебя в свой стремительный круговорот». «Мир — через стекло». 


Итак, с одной стороны, ТВ благо, с другой — таит опасность возникновения нежелательных качеств в личности ребенка. Но в этом противоречии нет ничего удивительного. Сами по себе изобретения и открытия—электричество, атомная энергия, законы генетики — могут служить на благо человечества или наносить вред. Судьба открытий зависит от того, в каких целях и в чьих интересах они используются. 

Также обстоит дело и с изобретением телевидения. 

Вот одно из подтверждений тому. 

Рост преступности, особенно среди молодежи, становится в наши дни бичом многих капиталистических стран. И целый ряд исследователей связывает это в какой-то мере с содержанием телевизионных передач. С этим нельзя не согласиться, если вспомнить тот шквал фильмов ужаса, насилий, убийств,  который ежедневно обрушивается с голубого экрана на  телезрителей США или Западной Европы. 

Возможно ли подобное в социалистическом обществе? Конечно же, нет. Но и в наших условиях телевидение, призванное стать большим другом человека, при неумелом пользовании им может принести вред. 

Давайте рассмотрим подробнее, какие функции присущи советскому телевидению и как к ТВ относятся взрослые. 

Вот данные ленинградских социологов. На вопрос: «Чем лично для Вас является телевидение?» наибольшее количество ответов (1265) гласило — средство развлечения, отдыха,  разрядки; 1156 человек — источник информации о событиях в стране и за рубежом. Для 899 зрителей телевидение —  средство расширения знаний и для 630 — средство приобщения к искусству. 

В глазах взрослого телезрителя, таким образом, ТВ  выступает в трех основных функциях: как средство отдыха, познания и информации. Эти функции помогают удовлетворить взрослым многие духовные потребности. Но поскольку сутки не  увеличились, то время, используемое на просмотр телепрограмм,  «вычитается» из времени, употребляемого на другие занятия. 

В этом отношении интересно привести данные из книги Б. Фирсова «Телевидение глазами социолога». На вопрос, что вы предпринимаете, чтобы посмотреть интересную передачу, на первом месте по количеству ответов оказались — «ложимся спать позже обычного». 

Если приходят гости, то вместо разговора с ними их  привлекают к просмотру телепрограмм; передачи заставляют  зрителей откладывать домашние дела, прекращать чтение книг, не ходить в день интересной передачи в магазины, прачечную, мастерские. 

Проходит время, к телевизору привыкают, но интерес к нему не ослабевает. Для большинства нет вечера без телевизора. Желание получить развлечение в домашних условиях со  временем перерастает в потребность быть постоянно в курсе  событий — политических, культурных, научных, спортивных. 

ТВ еще больше связывает рабочего, колхозника, студента, инженера, служащего, школьника с духовной жизнью всей страны. 

Стало нормой, когда знакомые люди, встречаясь друг с  другом на улице, в трамвае или в поезде, на работе или в гостях, начинают разговор не с «погоды», а с вопроса: «Вы вчера смотрели телепередачу?» И вслед за названием передачи  начинается обсуждение поднятых в ней проблем. 

Причем, если владельцы телевизора со стажем, то  разбираются не только близкие каждому общественные проблемы, но и художественные особенности спектаклей, концертов, фильмов да и любой телепередачи на вполне профессиональном уровне. Не раз автору этих строк приходилось слышать меткие  замечания о работе режиссеров, актеров, кино и телеоператоров, художников и музыкальных оформителей со стороны людей, стоящих далеко от искусства (как по профессии, так и по образованию). В адрес студии нередко приходят отзывы о  передачах, вполне удовлетворяющие требованиям, предъявляемым к профессиональным рецензиям и критическим разборам. И это не удивительно: телевидение ведет огромную эстетическую работу с населением. И. Андроников и С. Виноградова, А. Каплер и Б. Бабочкин, известные критики, журналисты,  писатели, художники — частые гости наших телезрителей; люди, создающие искусство, знакомят нас со своей творческой  лабораторией, со своими замыслами и профессиональными  «секретами». 

Телевидение как бы создает для каждого человека  постоянную сферу общения с интересными людьми на интересные темы, а продолжение этого общения происходит в кругу семьи, с соседями, с коллегами по работе. 

И, конечно, человек, в семье которого нет телевизора по «принципиальным» соображениям, порой нам кажется чудаком.

Ссылки на то, что многие передачи не могут удовлетворить взыскательного зрителя, мало убедительны. Нельзя же  отрицать огромную роль телевидения только на том основании, что встречаются передачи слабые или рассчитанные на «среднего» зрителя. 

Но разве статьи в газетах, журналах, а также книги люди читают подряд? И разве понятие «отбор» исключается, когда речь идет о телевидении? 

Вспоминается случай, который произошел лет 15 назад.  Центральная студия телевидения проводила довольно часто встречи со зрителями во Владимире, Ярославле, в Тушине, с работниками «Трехгорки»», рабочими завода «Манометр»» и т. д.

Вторую программу смотрело тогда весьма ограниченное количество зрителей. Из 3—4 часов, проведенных у  телевизора, иногда только час-полтора соответствовали интересам  данной семьи. Остальное время считалось потерянным на  ожидание, «а вдруг дальше будет что-нибудь и для нас интересное?». На встречи-дискуссии о программах телевидения  приглашались зрители, которые в письмах и устно выражали крайнее недовольство телепередачами. 

Режиссеры, редакторы шли на эти дискуссии вместе с  руководством телестудии робко, прячась за спины дикторов Леонтьевой, Кондратовой, Кириллова, Шиловой, которых, несмотря на все претензии к передачам, встречали как самых желанных близких знакомых. 

Начиналась дискуссия. Выступает одна группа людей с  резкими возражениями против обилия трансляций симфонической и классической музыки, длительных концертов, предлагая  взамен этого передать экран полностью хоккейным или футбольным состязаниям, современной эстраде, «модной» музыке и т. д. 

Но нам, телевизионщикам, не приходилось в этих случаях организовывать оборону. Прежде чем мы успевали объяснить  принципы формирования программ, необходимость учета   различных вкусов, другая группа зрителей, также пришедших  критиковать телевидение, яростно бросалась в спор, защищая классическую музыку, телепьесы, протестуя против  «безголосых», микрофонных певиц, против засилия второстепенных спортпрограмм и т. д. На наших глазах между зрителями  разыгрывались баталии. Мы едва успевали мирить участников дискуссии. 

Не переносятся ли подчас эти «телевизионные страсти» в семью? 

Правда, сейчас положение значительно изменилось. Зрители располагают несколькими программами, они  систематически через прессу, радио и телевидение получают  информацию о предстоящих передачах. Многие зрители уже по  названиям рубрик заранее определяют, что им следует выбрать для просмотра. 

Вырос уровень передач. При их планировании и подготовке учитываются интересы различных социальных групп зрителей. Телепрограмма стала более дифференцированной. 

Сейчас любой зритель, независимо от уровня культуры,  знаний, интересов, возраста, профессии, имеет возможность выбрать для себя то, что его интересует. 

Конечно, еще встречаются и слабые, и неудавшиеся  передачи, но это скорее исключение, чем правило. Может быть стоит относиться к передачам, как к ступенькам, по которым один человек прошел далеко, и поэтому многие передачи уже не представляют для него интереса, а другому еще предстоит  подниматься. При таком подходе — каждый должен найти свою «ступеньку», с которой он может продолжать восхождение к бесконечным высотам культуры. К сожалению, у некоторых ТВ превращается в своеобразное «хобби». И они становятся  телепленниками, проводя перед голубым экраном все свое свободное время, жертвуя своими остальными занятиями, порою и родительскими обязанностями. А родители, как известно,  отличаются от «неродителей» тем, что они должны все свои поступки обдумывать и оценивать с позиции пользы или вреда детям. 

Вот тут-то мы и подошли к одному очень важному вопросу. Кому все-таки принадлежит телевизор? Ведь вы, покупая  телевизор, вовсе не рассчитывали преподнести подарок своему ребенку — нечто вроде фотоаппарата, конструктора или  велосипеда? Вы утверждаете, что телевизор принадлежит вам и только вам, взрослым. 

Однако, как заметил один английский писатель, этим  «взрослым инструментом» очень быстро завладевают дети. 

Именно завладевают, т. е. быстро научаются им  пользоваться и смотрят передачи и со взрослыми, и без них, а в ряде семей дети становятся безраздельными его хозяевами и  потребителями. 

Это часто происходит даже в тех семьях, где родители  сознательно относятся к воспитанию детей и считают, что телевизор — важное воспитательное средство, которое отвлекает детей от улицы, от безделья, поучает, воспитывает, развивает, более того, дает ребенку то, что они, сами родители, дать не могут. В отдельных случаях родители относятся к телевизору как к своеобразной «теленяньке», эта позиция имеет свое  обоснование и ее можно в какой-то мере понять. Слишком уж различны условия жизни людей. 

Вот, например, что пишет гражданка Л.: «У нас в школе нет группы продленного дня. А во дворе много всяких хулиганов, которые слоняются по двору и от безделья озорничают. Я думаю, что сын больше пользы получит от телевизора, чем от таких ребят». 

Небезынтересно и другое письмо, также голосующее за «теленяньку». Это письмо пришло в редакцию из города  Львова: «Моей маленькой соседке Танечке пять с половиной лет. Ее родители пытались устроить в детский сад, но обычно она с первых дней там заболевала. Поэтому родители решили купить телевизор и оставлять ее дома одну. Она самостоятельно  включает телевизор и смотрит все подряд. Она раньше была  молчаливая, а теперь, когда родители приходят с работы, пытается рассказать то, что видела в передачах. Она стала петь, любить музыку. Родители серьезно подумывают о музыкальной школе. 

Другие родители из нашего дома рассказывают, как их дети такого же приблизительно возраста делают зарядку, рисуют и мастерят игрушки по заданию дикторов телевидения». 

Довольно убедительный пример, подтверждающий  положительные возможности телевидения, не правда ли? Телевизор взял на себя часть родительских обязанностей... 

У жителей отдаленных населенных пунктов свой довод: «Мы живем вдалеке от центров культуры, детям некуда пойти  развлечься, а телевидение компенсирует недостаток культурно-просветительных учреждений»». 

Однако довольно часто можно встретить родителей, которые вообще не задумываются о роли ТВ в воспитании детей, тем более, если детям 11—12 лет: «Уже взрослые! Пусть смотрят, плохого не показывают». Действительно, советское  телевидение в отличие от ТВ капиталистических стран не является коммерческим предприятием и не стремится путем  третьесортных передач щекотать нервы зрителей, не создает сенсации ради сенсации, не показывает сценки сомнительного в  нравственном отношении содержания. И все же: оправдана ли  такая переоценка ТВ? 

В связи с этим не могу не привести наблюдений, сделанных мною в течение месячного пребывания в одной из деревень Подмосковья. 

В трехкомнатной квартире живет семья Скворцовых: мать — медсестра, отец — столяр и дочь Люда — ученица 7-го класса. 

Девочка учится в школе во вторую смену, встает уже после ухода родителей на работу, завтракает и садится то ли делать уроки и одновременно смотреть утренние московские  телепередачи, то ли наоборот — смотреть передачи и между делом выполнять домашние задания. 
 
Можно сказать, что Люда учит уроки под аккомпанемент телевизора. Она читает или решает задачу, то и дело  поглядывая на экран телевизора, и если ей что-либо покажется  интересным, она усиливает громкость. Если передача музыкальная, она пытается одновременно заниматься, если же требует сосредоточенного внимания, Люда отодвигает тетради и только смотрит ТВ. 

Но вот Люда приходит из школы. Ужинают всей семьей вместе и одновременно у всех глаза устремлены на экран. Около десяти вечера родители укладываются спать. Подушки на кровати поставлены так, чтобы можно было сидя-лежа  смотреть телепередачи. Утомленные родители постепенно  засыпают, как-то вползая под одеяло. А дочь? Она сидит на диване и усердно продолжает смотреть передачи до окончания  программ. 

Я намеренно наблюдал семью, которая живет отнюдь не в стесненных жилищных условиях. 

Ведь подобные случаи, к сожалению, встречаются еще там, где у людей однокомнатная квартира, и родители, позволив себе в вечернее время посмотреть телефильм или концерт, невольно «включают» в состав зрителей и своих детей. Семья же, о которой я говорил, лишь подтверждает, что существует такая группа родителей, для которых вообще нет понятия режима или организации просмотра телепередач. 

Что же из себя представляет Люда? На первый взгляд ее можно отнести к числу «благополучных», скромных, спокойных детей. Но так ли это? Уже сейчас проявляется у нее душевная пассивность, безволие. Если приходится выбирать между  передачей или пионерским сбором, передачей или уроками,  передачей или книгой, она легко склоняется в пользу передачи. Люда любит спортивные и театральные передачи. Она ведет альбом, в который записывает данные о спортсменах и любимых  артистах. К сожалению, любовь к спорту, как и театру, этим и  ограничивается. Зимой она ни разу не ходила на лыжах, не каталась на коньках, на санках. По утверждению родителей, и в прошлом году телевизор заполнял весь ее досуг. 

О театре и говорить нечего. У девочки нет ни интереса к книгам о театре, ни желания критически анализировать  увиденные спектакли или участвовать в художественной  самодеятельности. 

Я познакомился с классом, в котором учится Люда.  Познакомился и с другими ребятами из этой школы — с четвертого по девятый; провел анкету, выясняющую интересы школьников, подсчитал, сколько времени уделяют они своему другу  «Голубому экрану». И убедился, что Люда не исключение, многие ребята 4—8-х классов посвящают ТВ от 2,5 до 4,5 часов  ежедневно. 

И, как правило, наиболее «активные» зрители не являются активными общественниками, не блещут знаниями по  школьным предметам. Не берусь утверждать, что в этом виновато только ТВ. Скорее вижу недоработку педагогического  коллектива, который мог бы и жизнь ребят сделать интересной да и родителям подсказать, как помочь этому. 

Но вернемся к Люде Скворцовой. Даже если допустить, что передачи, которые она просматривает, одна другой интереснее и каждая способствует ее духовному развитию, расширению кругозора, то все равно, как свидетельствуют специалисты, «поточный» метод просмотра передач — враг № 1 вашего ребенка. 

И начнем со здоровья. Сосредоточенный и длительный  просмотр телепередач вредно влияет на нервную систему, на  зрение ваших детей. 

Взрослые порой просто забывают, что возможности  детского организма ограничены, что он не приспособлен  выдерживать те же эмоциональные и умственные нагрузки, что и взрослые; даже познавательный процесс у детей проходит  интенсивнее, чем у взрослых. 

Вы часто говорите своим детям: «Не поднимай эту тяжесть, ты еще маленький, надорвешься». А разве долгое сидение у экрана телевизора не грозит им тем же самым? Правда, вы не сразу замечаете, как ребенок утомляется, как перенапрягается его нервная система, он становится раздражительным или вялым, плохо спит, не может сосредоточиться на уроках, быстро устает и т. д. 

Врачи-глазники и невропатологи категорически выступают за ограничение времени просмотра детьми телепередач. Вот что пишет врач-невропатолог Г. Долгопятов, отвечая на вопрос: 

«Ребенок у телевизора — опасно ли это?»: «В начале 50-х годов, в пору широкого внедрения  телевидения в наш быт, в медицинской литературе довольно часто стало встречаться понятие «телевизионная эпилепсия». Причиной этой эпилепсии является сильная эмоциональная нагрузка, повышенная чувствительность нервной системы к  мелькающему свету... И чем младше ребенок, тем скорее даже слабые раздражители (поток телесигналов, вызывающих реакцию) могут... привести к развитию неврозов». 

В зарубежной печати не раз отмечалось, что в тех случаях, когда детям разрешается без ограничения смотреть передачи, то у них наблюдается искривление челюсти и неправильный рост зубов. Специалисты связывают это с положением рук детей, смотрящих телепередачу: чаще всего они руками  подпирают подбородок. 

Не будем отрицать и даже брать под сомнения наблюдения западных педиатров за юными телеманами. Не будем  преувеличивать возможных последствий для нормального физического и умственного развития детей, в том числе и для зрения, но нельзя их и игнорировать: надо обеспечить ограничение,  разумную регламентацию просмотра передач детьми, стремиться к соблюдению телегигиены. 

Но здесь существует и другая крайность, о которой мы уже упоминали. Еще встречаются среди родителей ярые противники ТВ. Они считают его источником всех воспитательных бед, поскольку «телевидение отвлекает от приготовления уроков, нарушает режим дня, развивает созерцательность, уменьшает интерес к чтению, сокращает время прогулок и выполнения домашних хозяйственных дел». Бытует и такая точка зрения: «Телевизор вносит много дополнительных трудностей в воспитание ребенка, поэтому целесообразно его купить после  окончания сыном школы...» 

Как правило, так думают родители, заранее  предчувствующие, что им не справиться с конкуренцией нового  телевоспитателя, который полностью завладеет вниманием и временем их ребенка, сделает его послушным, ручным, безвольным  телесуществом и — как они считают — обеднит духовный мир их ребенка. 

«ТВ — это не первая необходимость, — так иногда  утверждают сторонники этой точки зрения, — Пушкин и Горький, многие крупные ученые и политические деятели прекрасно обходились без телеящика, значит, подлинная культура не в нем». 

Встречается еще один тип родителей, которые используют интерес детей к телевидению в качестве «кнута и пряника». За плохое поведение или учебу наказывают лишением права  смотреть телевизор и, наоборот, за хорошее выполнение домашних работ, послушание, получение высоких оценок награждают правом смотреть телевизор. 

Лишение просмотров телепрограмм за плохое поведение, двойки многим кажется вполне педагогичным. Ведь лишают же ребенка посещений кино, если он, допустим, начинает хуже учиться. Логика при этом проста: «ты тратишь время на кино, а учеба страдает. Поэтому будь добр, употреби это время на подготовку уроков». (Однако кино — это эпизод, а ТВ —  повседневность). 

Ради успешного учения ребенка его нередко освобождают даже от домашних, бытовых забот, создавая особые (я бы сказал неестественные) условия. Очевидно, таким путем  запрещения или разрешения можно добиться какого-то временного успеха «на одном фронте». Однако правилен ли этот путь? Думаю, что им можно воспользоваться только в виде  исключения. 

«Табу» на просмотр передач может быть применен как  единичный, конкретный акт, а не как метод  воспитания. Значительно важнее выработать у ребенка определенное, разумное избирательно-оценочное отношение ко всему, с чем сталкивается он в жизни. Телевидение — всего лишь частный случай. 

Некоторые сторонники ограничения, запретов, боясь вызвать обиду у детей, предлагают незаметно выключать  телевизор или вынимать предохранитель, а детям объяснять, что телевизор «устал», «спит»» и т. д. 

Спрашивается, а как долго можно, таким образом,  обманывать детей? Они же в 6—7 лет прекрасно управляются с  телевизором, а в 12—13 разбираются в принципах работы радио и телевизора. Юные техники под руководством старших сами собирают телеприемники. Так что спасительная ложь не спасет. 

			\ \\
И вообще надо стараться достигать воспитательного эффекта методами, которые со временем не опровергаются детьми и не разоблачают взрослых. 

Приведу довольно-таки яркий пример, иллюстрирующий, как далеко может завести родителей использование телевизора в качестве средства поощрения или наказания. 

Рассказала о нем одна мама, живущая в небольшом городке Одесской области. «В нашем городе не развеселишься:  кинотеатр, клуб, небольшой парк... Ну, ясное дело, для сына  единственное спасение — телевизор. Он готов смотреть передачи и днем, и ночью. Даже уроки не хочет делать, если идет  передача». 

Кстати, картина — многим знакомая: и тем, кто живет в небольшом городке или на окраине большого города —  «микрорайона» и тем более в деревне... Кроме того, разве в Москве, Ленинграде, Горьком и т. д. нет ребят, готовых уроки променять на развлечение, на телепередачи? И виноваты ли в этом небольшой город и телевизор? 

«Так вот, — продолжала мама, — мы решили интерес сына к телевизору использовать как стимул к учебе. Мы с отцом постановили: телепередачу .разрешать смотреть только в те дни, когда его будут спрашивать в школе или когда он сам будет активен на уроке. Во-вторых, за каждую тройку он получает право смотреть детскую получасовую передачу, за четверку — любую часовую, а за пятерку он становится владельцем  полуторачасового фильма или передачи»». 

— И помогает? — спросили мы. 

— А как же? Он активней поднимает руку, просит  учителей, чтобы его спросили. 

— А если он получит две тройки, то получает право на две получасовки? 

— Нет, на одну. Если он в один день получит и пятерку и тройку, то мы оцениваем его право смотреть передачу только по одной низшей оценке, т. е. по тройке. 

— Ну, а как же он в результате учится? 

— Вобщем средне, все бывает, но больше четверок, есть тройки»». 

Вы можете спросить, неужели эти родители всерьез  считают, что их сын потому и учится, что дома есть телевизор? Думаю, они верят в то, что нашли действенное средство как улучшить его учебу. Но не подумали о главном: о воспитании сознательного отношения к учению, о воспитании интереса к нему. Ведь только в этом случае может быть воспитано и правильное отношение к телевидению, и оно войдет в жизнь ребенка значительным событием, а не как приманка, награда или только развлечение. 

Кстати, с подобным отношением к телевидению как лишь средству развлечения приходится встречаться нередко. Многие так и считают, захотел отдохнуть — включай телевизор. 

Об этом говорит и специальное исследование, проведенное недавно в школах Москвы и Ленинграда. На вопрос, какое место в твоей жизни занимает телевидение, школьники разного возраста (и ученики начальной школы и выпускники) в  подавляющем большинстве (85—90\%) ответили, что оно помогает им отдохнуть и развлечься. 

Не поймите, что мы в какой-то мере хотим умалить роль  телевидения в удовлетворении потребности развлечься, отдохнуть. Более того, мы глубоко убеждены, что на телевидении еще мало создается веселых развлекательных передач, особенно для детей. Да, да, откровенно развлекательных, смешных до колик, где песни, шутки, приключения, комические ситуации  увлекательны, неожиданны, остры, свежи. 

Ведь не случайно, среди постоянных зрителей «Кабачка 13 стульев» (по существу единственной развлекательной  передачи, которая из месяца в месяц появляется на экране  телевидения) так много ребят самых разных возрастов. Это убедительно подтвердили опросы, проведенные в Вологодской области, в деревнях, городах и поселках Коми АССР, ленинградских и московских школах, Грузии и др. 

Но было бы абсолютно неправильно смотреть на телевизор только как на средство развлечения. Потому что телевизор может дать несравненно больше: стать источником знаний, политической информации, помощником в учебе, в воспитании детей и т. д. Но к этому мы еще вернемся. 

Давайте продолжим разговор об использовании телевизора как своего рода награды, приманки и т. п. 

Не кажется ли вам. что в подобном педагогическом приеме торжествуют формализм и равнодушие к сущности и  содержанию воспитания ребенка? 

И не воспитывают ли родители, таким образом, у этого  мальчика вместо интереса к знаниям всего лишь стремление во что бы то ни стало получить высший балл? 

Видится мне, что родители здесь любой ценой стремились добиться немедленных успехов, не думая о подлинных целях воспитания. 

Не могу не остановиться еще на одной категории родителей, объявляющих себя сторонниками теории «свободного  воспитания»». Чаще всего — это лишь способ оправдать свою  педагогическую бездеятельность. Употребленный термин очень мало имеет общего с известными в истории педагогики  концепциями Ж.-Ж. Руссо, Л. Толстого и других. В данном случае под  свободным воспитанием попросту подразумевается  предоставление ребенку полной свободы и, в частности, в просмотре  телепередач. 

			\ \\
Такие родители не утруждают себя целенаправленными и  скоординированными действиями и сводят свои родительские  обязанности лишь к заботе о здоровье ребенка, об удовлетворении его материальных потребностей. 

В одном из писем в редакцию читатель пишет: ««Пусть дети смотрят, что хотят и сколько хотят. Вы спросите, как мы  воспитывали наших двоих детей? Да никак не воспитывали: они сами воспитывались и оба выросли хорошими людьми». 

Что можно возразить автору этих строк? Во-первых, что  значит «наших детей никто не воспитывал, они воспитывались сами»? Они ведь росли не в изоляции. И если даже не  посещали детский сад, то учились в школе, где воспитанием  занимаются учителя, пионерская и комсомольская организации. А можно ли отрицать воспитательное воздействие участия в работе школьных кружков, в различных школьных  начинаниях: субботниках, вечерах и т. д.? 

Во-вторых, разве дети были изолированы от детской  литературы, кино и, наконец, радио и телевидения? Вряд ли в наших условиях можно сказать о детях, что их никто не воспитывал. 

Можно только сожалеть, что в этом большом  воспитательном оркестре очень важная партия семейного воспитания  звучала, как можно понять из письма, очень тихо. 

В семье инженера К., наоборот, с первых дней после  покупки телевизора было решено создать традиции: ежедневно,  управившись с домашними делами, всей семьей садиться в 8  часов вечера за стол ужинать и одновременно смотреть ТВ,  сначала программу «Время», а затем кинофильм, спектакль или музыкальную программу. В семье К. двое детей — учащиеся 7-го и 9-го классов. Вместе с родителями они часто сидят у экрана до окончания вечерних передач. 

Расчет родителей прост: они хотят, чтобы по вечерам  создавалась хорошая семейная атмосфера. Детям, однако,  разрешается включать телевизор и днем. 

В конечном итоге получается, что ребята в этой семье почти каждый день просиживают у телевизора от двух с половиной до пяти часов. 

И что же? Как будто довольны и взрослые, и дети.  Телевидение вполне оправдывает ожидания: оно развлекает,  информирует, является источником знаний и т. д. Семья стала  дружней, взаимопонимание между младшим и старшим поколением явно выросло, но... В семье инженера К. стали меньше читать, реже ходить в гости, в кино, ездить за город кататься на лыжах и т. д. Дети предпочитают теперь «болеть» у телевизора, и сами стали реже ««гонять» мяч во дворе со своими  товарищами. 

Исследования, проводимые в нашей стране и за рубежом, говорят о том, что дети из «телевизионных» семей по  сравнению с детьми из семей, где телевизора нет, меньше играют, рисуют, лепят, мастерят. 

Но можно ли ставить этот факт в прямую зависимость только от наличия или отсутствия телевизора? Связано ли это  «меньше» или «больше» только с той пассивностью, которая, как часто любят говорить, развивается у постоянных телезрителей? 

Во-первых, дети по своей натуре и без телевидения делятся на пассивных и активных. Если придерживаться этого наиболее простого и понятного деления, то пассивные дети при  отсутствии руководства со стороны взрослых, под влиянием  телевидения становятся еще более пассивными («телевизионные дети»). 

На активных детей, как показывают наблюдения,  телевидение действует благотворно. У них появляется желание  попробовать свои силы в той или иной области деятельности. Как правило, активный темперамент не делает их рабами экрана или пассивными созерцателями жизни через стекло.

Итак, пассивные становятся еще более пассивными... Можно ли бороться с этим явлением? Да, мы решительно хотим подчеркнуть, что этот вывод верен лишь постольку, поскольку не учитывается одно важное обстоятельство, а именно, что дети смотрят передачи не в одиночку. Как правило, ребенок не  находится с телевизором «с глазу на глаз». Многие передачи юные зрители смотрят в кругу семьи, вместе с родителями, более того, на их воспитание огромное влияние оказывают взрослые. Поэтому весьма важно — какое это влияние, как сами взрослые относятся к передаче, стремятся ли они превратить ТВ в своего союзника в воспитании ребенка. А в первую очередь это значит пытаются ли они пассивность преодолевать, активность  развивать? Вы спросите, как это сделать? Каких надо  придерживаться правил? Как на практике превратить ТВ в полноправного воспитателя ребенка? 

Это, безусловно, сложный процесс, во многом  индивидуальный, своеобразный для каждого случая. Поэтому родителям в первую очередь надо самим серьезно поразмыслить над этими вопросами. Мы же со своей стороны остановимся на  некоторых общих правилах и принципах, которые в  значительной мере подсказаны наблюдениями, практикой. 

Прежде всего отметим, что телевизор не вносит чего-либо нового в педагогику, но, несомненно, обогащает приемы  воспитания. Телевизор заостряет ряд проблем и прежде всего  требует более осознанного, активного и ответственного отношения к воспитанию детей. 

Появление в семье телевизора создает ситуацию, когда  воспитанные ранее привычки и навыки подвергаются испытанию. 

Например, обычно ребенок, придя из школы, обедает, немного гуляет и садится за домашние задания. Он привык это делать ежедневно. Но вот теперь, вернувшись из школы, ваши сын или дочь встречаются с немалым соблазном: спортивная телепередача и... прогулка откладывается. 

После спортивной передачи объявляется кинофильм или другая интересная передача, и... откладываются уроки. Ведь их можно сделать и позже, а вот фильм (или передачу), возможно, в другой раз и не увидишь. 

Вот вам и первое испытание, когда подвергается проверке, насколько прочно укоренилась у вашего ребенка привычка выбирать главное, соблюдать ранее заведенный распорядок. 

Другой пример. С появлением в семье телевизора некоторые родители считают, что теперь они могут позволить себе иногда «отдохнуть» от резвого ребенка и получить «минуту покоя»»,  усадив его перед экраном телевизора. 

Таким образом, сокращается время общения ребенка со взрослым, которое так необходимо ему для полноценного  развития. 

Телевизор нисколько не освобождает родителей от общения с детьми, наоборот, он создает дополнительную возможность поговорить на тему увиденного на экране. 

Мы с вами уже достаточно уделили внимания отношению детей и взросых к телевидению. Теперь попытаемся понять, в чем «же заключаются для них магнетические свойства ТВ? Почему дети чуть ли не с трех лет становятся добровольными пленниками голубого экрана? В чем тут секрет? 

Вряд ли правильно будет ссылаться на «безволие» или  «неразумность» детей, или на ««сверхъестественные» свойства ТВ. Суть в психологии детского возраста, в особенностях развития духовного мира детей, в психологии их восприятия. 

Такие качества ребенка, как любознательность, стремление к познанию окружающей действительности, являются  естественными и необходимыми для нормального существования человека. 

Отсюда и тяга детей к экрану телевизора, который,  воздействуя на органы зрения и слуха, предоставляет ребенку  широкую информацию о внешнем мире, помогая ему  ориентироваться в явлениях природы, событиях общественной жизни. Народной мудростью прекрасно выражена одна из  особенностей человеческого восприятия: ««Лучше один раз увидеть, чем сто раз услышать». В своей «Книге о живописи», размышляя над этим свойством человеческой психики, Леонардо да Винчи писал: «Глаз, называемый окном души, — это главный путь, которым общее чувство может в наибольшем богатстве и  великолепии рассматривать бесконечные творения природы, а ухо является вторым». 

Телевидение же воздействует как на первое «окно души», так и на второе, хотя в основе — все же изображение. Более того: на экране перед нами происходит постоянная смена  изображения, имитация движения (как говорят специалисты,  телевидение владеет и кинетическим образом). Это особенно притягательно для детей. Как и театр, кинематограф, ТВ — явление синтетическое, т. е. оно владеет выразительными средствами других видов искусства: оно соединяет в себе возможности театра, кино, радио, документальной хроники, литературы и др. Оно объединяет самые разнообразные жанры. В самом деле, ТВ, как говорят, «всеядно»». Оно способно передать короткую, но впечатляющую информацию о событии, репортаж из  космоса или о спортсоревновании, показать документальный фильм о Галапагосских островах и передать интервью с  ученым, сидящим в студии; опера и драма, балет и цирк  соседствуют на экране с научной викториной, учебная передача — с кукольным представлением, художественный фильм — с  лекцией и т. п. 


Впрочем, перечислить все, что может телевидение, было бы, я думаю, невыполнимой задачей. Причем сила воздействия не только в наглядности зрительного образа, лежащего в основе синтетического языка телевидения, но и в использовании  возможностей музыки, слова и др. Не случайно, что любая  передача — научная ли, информационная, спортивная — плод совместной работы сценариста, редактора, режиссера,  художника, звукорежиссера, звукооформителя, оператора и др. Даже если речь идет о репортаже с места событий, зритель не просто «видит» событие, он видит его как бы глазами режиссера и оператора. Они то укрупняют его, то выделяют детали, придают определенный ритм передаче (напряженный, радостный,  тревожный и т. д.), дают возможность телезрителю сопоставить свое впечатление с реакцией зрителей, сидящих в зале. Все это усиливает восприятие происходящего. Если необходимо, дают замедленную съемку, вводят мультипликацию (например, при показе невидимых глазом скрытых процессов, происходящих в клетке, или при химических реакциях), используют повтор кадра в футболе, хоккее и др. 

Но ведь все это можно увидеть и в кино и вообще, наверное, у вас давно назрел вопрос: а в чем собственно отличие  телевидения от других средств массовой коммуникации? Ведь,  конечно, не только в том, что телезрителю почти не требуется никаких усилий — щелкнул выключатель и акт познания начался. 

Вот здесь, очевидно, самое место коснуться вопроса — ТВ в кругу собратьев. 

Телевидение называют младшим братом кино. На самом деле — не только кино. У ТВ имеются ««родственные»» связи и с другими видами массового общения — периодической печатью, радио,театром. 

Рассмотрим подробнее, что между ними общего и в чем их отличие? Газета и журнал, радио и ТВ рассчитаны на общение с миллионами людей, но происходит оно в условиях семьи, маленькой камерной аудитории. Отсюда — своеобразие  художественных средств выражения. 

У радио и ТВ — одна аудитория, одинаковые условия  восприятия, содержание радиовещательной и телевизионной  программы во многом сходны. Такие функции радио, как передача информации, распространение знаний, эстетическое  воспитание и т. п., совпадают с функциями ТВ. Поэтому радио и ТВ используют жанры, которые ни в кино, ни в театре не встретишь (беседа, концерт-загадка, ««клубные»» передачи,  международные обозрения, ««панорамы»» и т. д.). 

Однако ТВ не может и не должно заменить радио в  воспитании детей. 


Отсутствие у радио такого качества, как «зрелищность», 
создает, как это не парадоксально звучит, определенные  
преимущества. Например, оперируя только звуком, радио способно 
сильнее влиять на воображение, развитие фантазии,  
абстрактного мышления. 
Радиопередачи не требуют абсолютного сосредоточения, 
как требует того телевизор. Во время радиопередач человек 
может одновременно выполнять какую- либо работу, что и 
делают многие. 
Хорошая симфоническая музыка, например,  
воспринимается по радио гораздо глубже, так как зрительные образы — 
оркестр, участники хора, музыкальные инструменты и т. д. — не 
31 
отвлекают радиослушателя от сосредоточенного и  
полноценного восприятия музыкального произведения. 
Телевидение требует значительной мобилизации  
психофизических сил, а поэтому больше утомляет человека, вызывая 
сильное напряжение нервной системы. 
Однако телевидение обладает своими преимуществами — с 
помощью звукозрительных динамичных образов оно более 
доступно и понятно людям, оно дает возможность не только 
услышать, но и увидеть балет и скульптуру, далекие страны и 
«подлинных»» людей. 
Зрительная информация, отражая конкретные образы  
окружающей жизни, быстрее и основательнее воздействует на 
сознание человека. 
Зная преимущества и слабые стороны радио и ТВ, родители 
должны осознанно выбирать передачи с учетом конкретных 
воспитательных задач. 
У ТВ есть еще один близкий родственник — театр. 
У них одна общая особенность, которая очень  
привлекает зрителей, активизирует их восприятие. Телезрители и  
зрители в театре видят происходящее событие в то самое  
время, когда оно происходит (мы, естественно, исключаем  
случай записи событий на кино- или магнитную пленку). 
Но в отличие от телезрителя, который вынужден  
довольствоваться лишь оптическим эффектом, бледной тенью, зритель в 
театре видит перед собой живых исполнителей, во плоти и 
крови, физически ощущает их близость, их «дыхание»». Если 
представить только на секунду сцену из балета, которой мы 
наслаждаемся в театре во всем ее физическом блеске и  
великолепии, во всем ее человеческом масштабе в союзе с мощью 
оркестра, пением скрипок и звуков литавр, и ту бедную копию 
балета, которую мы видим на экране ТВ, то сразу становится 
ясным, что ТВ не соперник театру, у каждого свое  
преимущество, у каждого свое лицо. И никакой цветной экран никогда не 
заменит ни театрального зрелища, ни театрального  
праздничного настроения. 
Или еще одно: зритель в театре всегда.ощущает дыхание 
32 
зала, заражается общим настроением, он вместе со всеми  
воздействует на игру актеров, подбадривает их, настраивает их, 
общается с ними, становится участником спектакля. 
Телезритель лишен возможности воздействия на актера. 
Если театральный зритель активен и выбирает из  
происходящего на сцене то, что его больше всего интересует в данный 
момент, то телезритель в этом смысле пассивен, он видит 
только то, что ему решил показать режиссер — деталь или 
общий план, актера или... публику в зрительном зале. 
Однако телезритель имеет свои преимущества — крупный 
план актера (или выступающего) дает возможность ему с 
большей психологической глубиной проникнуть в образ 
героя, увидеть его переживания, глаза, мимику. 
Отдельные детали, показанные крупным планом или в 
определенном ракурсе, могут служить очень важным средством 
раскрытия идеи спектакля или передачи. Конечно, его  
впечатления ни в коей мере нельзя сравнить с тем, что видит 
театральный зритель, сидящий в последнем ряду или на  
балконе. Они разные, несравнимые. Вот почему можно с  
уверенностью сказать, что развитие ТВ нисколько не мешает развитию 
театра. И сейчас на интересные театральные спектакли  
невозможно достать билета. 
Подлинная «театральная культура» человека не может  
развиваться полноценно только благодаря регулярному просмотру 
телеспектаклей по телевидению. 
Дорогие родители! Если у вас есть малейшая возможность 
приохотить вашего ребенка к театру — не упускайте ее!  
Воспитание театром приносит человеку радость и заметно, от  
спектакля к спектаклю, обогащает его высокими идеями и чувствами, 
как говорил в свое время К. С. Станиславский, творит  
человеческий дух. 
Очевидно, ТВ ближе всего по восприятию стоит к  
кинематографу. 
У них действительно очень много общего. Прежде всего, их 
роднит экран —двухмерное световое изображение, достоверно 
отображающее жизнь так; как мы ее видим на самом деле. И 
33 
кино, и телевидение обладают многими, своеобразными  
выразительными средствами. 
Но если говорить о том, чем они отличаются, то кино — 
всегда «вчера»», ТВ — всегда ««сегодня»». «Эффект присутствия»» 
— это огромное преимущество ТВ. 
Когда-то на заре создания первых фильмов публика  
шарахалась от двигающегося на экране поезда, принимая  
изображение за действительность. 
Сейчас уже никто не воспринимает фильм как сиюминутно 
происходящий факт. 
Отсюда и совершенно иной характер восприятия. С  
помощью датчиков изучалось состояние заядлых футбольных 
болельщиков, присутствовавших на стадионе и одновременно у 
телевизора. Показатели были довольно близкими. Для  
сравнения взяли еще одну группу таких же рьяных болельщиков, 
которые не видели с£утбольную игру и не знали, как она  
проходила, не знали они и результатов. На другой день им полностью 
показывался фильм, который подробно фиксировал всю игру. 
При сравнении оказалось, что интенсивность переживания у 
киноболельщиков была значительно ниже, чему у  
телеболельщиков: игра для них хоть и была интересна, но это был все же 
«вчерашний»» день, своего рода история, а не живая жизнь. 
В телевидении трудно «сжать»» время, как это делает кино 
Например, ученик подходит к школе, открывает дверь. Далее 
режиссер показывает, как этот ученик поднимается по лестни 
це, подходит к классу, открывает дверь, садится за парту. Такик 
образом, можно проследить весь путь ученика от школьное 
двери до класса и т. д. В телевидении «телевремя» равне 
реальному времени. 
Кино способно не только сжимать, но и растягивать время. 
Фильм, снимаемый несколько месяцев, может менять десятки 
мест съемок. 
Телепередача, телеспектакль всегда ограничены, они хотя и 
используют больше мест действий, чем театр, и все-таки в 
десятки раз меньше, чем кинем-атограф. 
Поэтому телеспектакли избегают массовых сцен и обилия 
34 
мест действия. Все это ограничивает возможности ТВ, но оно 
пытается компенсировать их за счет «крупных планов»,  
психологических сцен, раскрытия внутреннего мира героев. 
Конечно, такой фильм, как «Чапаев», выдающееся  
произведение советского киноискусства, в котором большое место 
занимают очень точные и удивительно выразительные  
массовые сцены, не следует смотреть по ТВ. Так, например, широко 
известная по своей силе воздействия сцена «психической  
атаки» на экране даже самых больших телевизоров выглядит 
совершенно иначе, чем в кино. А чего стоит дыхание  
кинозала... Смех в зале! Аплодисменты, душевный подъем, который 
создается а результате сопереживания с сотнями таких же 
зрителей. 
Но, наверное, вам приходилось наблюдать случаи, когда  
некоторые фильмы у одних людей вызывают грусть, у других — 
смех, у одних — веселое настроение, у других — возмущение. 
Здесь все зависит от настроя, душевного состояния зрителя и 
склада характера. Очевидно, такой фильм лучше смотреть по 
телевизору, где не будет столкновений столь бурных, разных 
настроений, мешающих восприятию. 
Известный французский кинорежиссер Анри-Жорж Клузо 
говорит, что телевидение может показать глаза Робеспьера, 
руку Робеспьера, но оно не покажет революции. Кино  
позволяет увидеть революцию, толпу, но оно в то же время может и 
сосредоточить внимание на отдельных деталях, предметах  
одежды и т. п. Родители, зная законы большого и малого экрана, 
законы восприятия в кинозале и в условиях семьи, должны 
понимать, что телевидение — это не домашнее кино и оно ни в 
коем случае не заменит посещение кинотеатра, его  
воспитательных возможностей и той радости, которую приносит  
хорошая картина, просмотренная в зале вместе с другими ребятами. 
Телевидение невидимыми нитями связано еще с одним 
видом массовой коммуникации — печатью. 
Пресса, как и радио, и ТВ, вездесуща, она проникает в 
каждый дом, предприятие, школу. Она также многожанрова и 
многотемна. Но сравните ваши впечатления, когда вы, допу- 
35 
стим, одну и ту же речь читаете в газете, слышите по радио или 
воспринимаете по телевидению. 
На телеэкране перед вами — говорящий человек. У него 
естественная речь. Она, может быть, порою и лишена той  
выверенное™, отточенности каждой фразы, каждого слова, что 
присуще печати. Говорящий может и повторяться, быть не  
всегда логичным, даже сбиваться. Но вы присутствуете при  
рождении мысли. Да и содержание воспринимаете через интонацию 
говорящего, его жестикуляцию и мимику, выражение глаз, 
наконец внешний облик и тембр голоса. То есть телевидение 
в сравнении с печатью обладает множеством дополнительных 
средств воздействия и раскрытия содержания. 
В этом огромное преимущество ТВ. 
По мере развития телевидения оно все больше и больше 
будет обогащать свои выразительные средства за счет других 
средств массовой коммуникации и видов искусства. Но вряд ли 
правомерно сравнивать телевидение с его собратьями в семье 
искусств в плане «лучше или хуже». У ТВ есть своя собственная 
сфера, свое собственное достойное место под солнцем, свои 
задачи, свои высоты. 
Правда, в капиталистических странах не стихают голоса, 
пророчащие, что телевидение подменит традиционные формы 
культуры и приведет к стандартизации вкусов, формированию 
усредненных культурных потребностей, утверждению  
эстетических критериев, не имеющих ничего общего с действительными 
ценностями. 
Мы, конечно, не можем согласиться с тем, что пороки  
буржуазного ТВ, пропагандирующего «массовую культуру»,  
распространяются на природу самого телевидения. В  
капиталистическом обществе господствующие классы заинтересованы в  
создании культуры для элиты и культуры для масс, которая 
стандартизирует и нивелирует вкусы, устанавливает  
духовную моду, угождая непритязательным мещанским требованиям. 
С помощью довольно высокой техники и технологии  
телевидение способствует формированию бездумного пассивного  
обывателя, поглощающего телеудовольствия, различного рода раз- 
36 
элекательную продукцию и сенсационнообразную  
информацию. ««Массовая культура» погружает телезрителя в  
иллюзорный мир, отвлекая его от острых социально- экономических 
проблем действительности. 
Массовая культура — это изобилие комиксов и детективов, 
фильмов ужаса и слащавых назидательных историй о том, как 
хорошо живется хорошим рабочим в современном обществе 
«народного» капитализма. Это феерические шоу,  
сдобренные стриптизом, это сенсации-однодневки, которые  
подменяют информацию и т. д. И все это «цементируется»  
бесконечной рекламой товаров, различного рода буржуазной  
пропагандой. На привлекательность «упаковки» этой духовной 
пищи, на рекламу ее капиталистическое государство не  
скупится! 
Армия социологов изучает вкусы «среднего» гражданина, 
чутко следит за изменчивостью моды, помогая сохранять  
необходимый уровень действенности изделий «массовой культуры». 
В Советском Союзе и социалистических странах не  
существует самого понятия «массовая культура». В основе развития 
социалистической культуры лежат положения, высказанные 
В. И. Лениным в его знаменитой беседе с.Кларой Цеткин: 
«... Искусство принадлежит народу. Оно должно уходить  
своими глубочайшими корнями в самую толщу широких  
трудящихся масс. Оно должно быть понятно этим массам и любимо 
ими. Оно должно объединять чувство, мысль и волю этих масс, 
подымать их. Оно должно пробуждать в них художников и 
разбивать их». 
То есть речь идет о громадной работе, направленной на 
подъем культурного уровня всего народа, чтобы то, что  
составляет духовную ценность человечества, стало достоянием 
каждого советского человека. 
Советское телевидение, являясь массовым средством  
распространения культуры, стремится эстетически развивать  
зрителя, повышать его духовный и художественный уровень и 
потребность в общении с подлинным искусством. 
Итак, мы с вами установили, что существенным признаком 
телевидения является поток разноообразной информации — 
37 
художественной и документальной, краткой и длительной, 
емкой, разножанровой и многотемной, способной  
удовлетворить различные интересы, вкусы, познавательные потребности 
человека. Телевидение в силу своей динамичности, зрелищно- 
сти, наглядно-звуковой природы очень полно, выразительно и 
доступно отражает объективный мир. Более того, поток  
передач, взаимодополняющих друг друга, отражающих разные  
стороны жизни, в значительной мере фиксирующих «сегодняшний» 
день, текущие события, как бы воссоздает атмосферу «живого 
дыхания», общественный ритм и среду, в которой живет и 
действует человек. Более того, иногда экран телевизора 
создает иллюзию личного присутствия при том или ином  
событии. 
Мы подошли, таким образом, еще к одному магнетическому 
свойству телевидения, о котором уже упоминали выше,  
называемому в специальной литературе ««эффектом присутствия». Он 
основан на показе, трансляции события в момент его  
свершения (этот эффект также называют «сиюминутностью»). 
Телевидение может показывать событие непосредственное 
момент его свершения, доставляя, так сказать, информацию с 
места действия прямо на дом. Такая информация обладает 
особой психологической притягательностью. Этот эффект 
сиюминутности буржуазное телевидение часто создает 
искусственно, чтобы привлечь телезрителей к экрану. 
Созданию эффекта сиюминутности способствует и стиль 
работы дикторов и выступающих перед телеэкраном. 
Особенно он действует на детей. Наверное, вы сами не раз 
бывали свидетелями, как они здороваются и прощаются с  
диктором, отвечают вслух на его вопросы. 
Была однажды передача, в которой Кот-книгоноша  
(персонаж кукольной сказки), распевая песенку про Машу-растеря- 
шу, грозился приехать к тем ребятам, которые не хотят сами без 
помощи взрослых обуваться и одеваться. И, как потом  
сообщали родители, дети, которые действительно заставляли мам 
или бабушек одевать себя, в тот момент, когда кот садился на 
38 
велосипед и отправлялся в путь, прятались под стол, под 
шкаф,чтобы Кот не нашел их. 
В другой раз Кот-книгоноша обещал заехать к тем ребятам, 
которые всегда капризничают за столом, не хотят пользоваться 
ложкой и вилкой. Из писем редакция узнала, что многие дети, 
хотя их и «миновала кара Кота-книгоноши», стали неузнаваемы, 
исправились. 
Когда ребенок мал, для него самый главный авторитет — 
мама и бабушка, папа или дедушка. Когда ребенок становится 
школьником, то у него появляется новый авторитет —  
учительница. Она главный судья во всем. И пусть папа — доктор  
педагогических наук, а учительница, Марья Ивановна, работает в 
школе после педучилища первый год, все равно ребенок может 
со слезами на глазах возражать: «Нет, папе, ты не прав, Марья 
Ивановна сказала, что синий карандаш лучше...»» 
Но вот в семье появился телевизор. Школьник, однажды 
придя в класс и набравшись храбрости, встает и говорит: 
«Марья Ивановна! А вчера по телевизору показывали не так, 
как Вы говорите. Я сам видел собственными глазами»». Для 
этого ученика телевизор стал главным арбитром, судьей. 
Телевидение обладает высшим авторитетом (особенно это 
относится к школьникам среднего и старшего возраста) не 
только потому, что оно представляет общество в целом, но и 
потому, что позволяет самим ребятам видеть и самим судить об 
увиденном. 
Таким образом, собственное суждение, подтверждаемое 
виденным своими глазами, обретает большую моральную силу. 
А как же быть теперь с родительским авторитетом? Он стал 
слабее? Родителям теперь стало труднее? «Да»»,— скажем 
мы,— в том случае, если родители будут игнорировать этот 
«сдвиг авторитетов»». «Нет»,— скажем мы, — если родители 
отнесутся со всей серьезностью к новому воспитателю и будут 
учитывать его возможности и роль в воспитании детей». 
ТВ может и должно способствовать укреплению  
внутрисемейных связей и авторитета родителей, но произойдет это не 
автоматически, а при определенных усилиях со стороны самих 
39 
родителей. Каких усилий? Давайте остановимся, поразмыслим, 
попробуем ответить на следующие вопросы: 
1. Помогает ли ТВ укреплению взаимопонимания и  
взаимоотношений между членами вашей семьи? Способствует ли ТВ 
утверждению, укреплению вашего родительского авторитета? 
Если «да», то почему именно? 
Если «нет» —тоже почему? 
2. Какие сознательные и систематические шаги вы  
предпринимаете, чтобы ТВ не только расширял кругозор вашего  
ребенка, но и влиял на формирование его характера, привычек,  
взглядов'? 
Нам кажется, что раздумья по поводу предложенных  
вопросов дадут вам правильные ориентиры. 
А сейчас хотелось бы остановиться еще на одной  
характерной особенности ТВ, привлекающей детей к экрану. Это скорее 
особенность не ТВ, а отношений между ТВ и юным зрителем. 
Это добровольность. Добровольность в решении ребят  
смотреть телепередачи. Ребенок в семье или в школе постоянно 
является «объектом воспитания». Его без конца воспитывают 
словом, тоном, примером, требованием, наказанием или 
поощрением. 
И даже тогда, когда давно' миновал тот возраст, когда 
ребенок у всех и по всякому поводу спрашивал: «А почему?», и 
он уже давно старается постигать жизнь самостоятельно,  
родители часто по-прежнему продолжают поучать, опекать,  
навязывать свои представления, предъявлять бесконечные  
требования, слабо аргументированные. Ребенок постоянно испытывает 
прямое и открытое педагогическое воздействие со стороны 
взрослых в школе и дома. 
А ТВ? Смотреть ТВ ребят никто не обязывает. Это  
добровольное дело. И приятное. И неподотчетное. 
Юный зритель усаживается к экрану по своему  
собственному желанию. Он сам выбирает передачу. Ее усвоение никем 
не контролируется, оценки за это не ставятся и т. д. 
У школьника создается полная иллюзия, что вот уж тут-то его 
никто не воспитывает, он сам смотрит то, что ему хочется и что 
ему интересно. 
40 
На самом деле, как вы прекрасно понимаете, каждая  
передача д детей выполняет свою определенную педагогическую 
функцию. В отличие от киноаудитории телеаудиторию можно 
считать постоянной, а общение ТВ с юным зрителем —  
систематическим. И следовательно, воспитательный  
процесс—непрерывным и целостным. 
Как мы уже говорили, телевидение привлекает к себе юных 
зрителей и стилем своего отношения к нему. Заметьте, дикторы 
и выступающие по телевидению не назидают, не поучают, они 
предлагают, советуют или даже сами спрашивают совета у  
зрителя, часто оставляют вопрос открытым (««А как вы думаете, 
ребята?»»), чтобы дети сами задумывались над ответом. К  
телезрителю обращаются в доверительном, уважительном тоне, 
диктор как бы апеллирует лично к каждому зрителю, тем более, 
что техника ТВ дает возможность выступающему смотреть 
прямо в глаза, усиливая чувство личной связи, соучастия и т.д. 
Конечно, стиль взаимоотношений, основанный на подлинном 
авторитете взрослого, не привилегия только ТВ, он может быть 
использован и в семье, и в школе. Кстати, чаще всего он  
устанавливается в тех семьях и с теми учителями, которые не 
придерживаются в своей воспитательной практике так  
называемой «авторитарной педагогики»», т.е. не пользуются грубо 
своим правом, властью, силой. 
А чтобы наглядно продемонстрировать тезис о  
добровольности, как притягательной силе ТВ, попробуем представить себе 
ситуацию: родители и учителя по своему усмотрению без учета 
интересов и желаний школьников составляют расписание 
домашних просмотров телевизионных передач. 
Мало того, осуществляют контроль в полном смысле этого 
слова за усвоением просмотренных передач, ставят оценки, 
журят за плохие, поощряют за хорошие, наказывают за пропуск 
просмотра передач, как за пропущенный урок... 
И тогда, я в этом нисколько не сомневаюсь, дети  
немедленно потребуют весенних, зимних, летних и даже осенних 
каникул... от телепередач. 
41 
Итак, дорогие родители, мы надеемся, что нам удалось 
выполнить первую задачу, а именно: показать, чем является ТВ 
для детей и взрослых, каковы его качества и возможности с 
точки зрения воспитания подрастающего поколения. Очевидно, 
читая нашу книжку, вы вспоминали ситуации, подобные  
упомянутым, из своей жизни и дополняли их личными впечатлениями. 
Благами телевидения человек должен пользоваться  
сознательно. При неправильном использовании ТВ, этого современного 
технического чуда, оно может принести и значительный вред. 
Таким образом, при пользовании телевизором следует  
исходить из двух условий. 
Условие первое. Это понимание той простой истины, что 
появившийся в семье телевизор не может быть нейтральным. 
Телевизор — это не предмет комфорта и не только  
свидетельство благосостояния. Как только зажегся огонек  
телеприемника, полился его голубоватый свет — ваши ум и чувства  
включаются в сложный процесс восприятия увиденного и  
услышанного. Вспоминайте всегда в этот момент о детях. Ни одна передача 
не пройдет мимо их пытливого взгляда, живого воображения. Я 
позволю привести здесь слова работника чешского  
телевидения Зд. Михальца, который считает, что, когда взрослые, желая 
«успокоить» своих детей, сажают их к экрану телевизора, не 
задумываясь, а стоит ли им смотреть данную передачу, то это 
все равно, что выпроводить ребенка во двор, на улицу, не зная 
в какой компании, хорошей или плохой, он там окажется. 
Условие второе. Суть его очень проста. Если вы не на 
словах, а на деле серьезно относитесь к воспитанию ребенка, 
не на словах, на деле, любите его, будьте всегда готовы  
ограничить свои собственные желания и в просмотрах телепередач. 
В некоторых семьях с первых дней покупки телевизора  
установилась традиция усаживаться после работы к экрану ТВ и 
смотреть все подряд до конца передач. По вечерам тогда нет 
места полезным домашним делам или любимым занятиям, будь 
то игра в шахматы, вязание, чтение или рисование. 
Приходится, к сожалению, наблюдать, как в некоторых 
семьях это перерастает в привычку проведения свободного 
42 
времени и становится почти единственным его способом. Кто- 
то очень метко сказал, что в наше время ТВ — это массовое 
хобби. 
Иные родители недопонимают, что учеба для ребенка — это 
серьезный труд. Он требует тишины, сосредоточенности, 
нервного напряжения. Всякие посторонние помехи не только 
затрудняют ребенку его работу, но и значительно утомляют его. 
Поэтому, наше второе условие — сознательно ограничивать 
свои желания, в том числе просмотры телепередач, в интересах 
ребенка, если они: 
а) могут отвлекать его внимание от выполнения домашних 
заданий, 
б) отвлекать от полезных игр или других занятий, увлечений 
(например, от чтения), 
в) если ребенок под влиянием взрослых передач раньше, 
чем это диктуется его возрастом, начнет вникать в проблемы 
взрослой жизни. 
Студии телевидения, как правило, при составлении  
программы на день, сугубо взрослые передачи стараются ставить 
в позднее время, т. е. после 21 часа 30 минут. 
Но об этом подробнее в третьей главе, специально  
посвященной программам. 
Теперь об азбучных истинах или о советах, которым будет 
посвящена следующая глава. Они доступны и просты в  
«употреблении» для любой семьи. Так же, как и буквы при сложении 
слов, они всего лишь основа для решения конкретных  
воспитательных задач, которые встают в семьях, где есть дети и  
телевизор. 
43 
ГЛАВА С* 
АЗБУЧНЫЕ ИСТИНЫ 
Азбучная истина № 1, или 
«Делу — время, потехе — час»». 
Из своего личного опыта, знакомства с письмами  
телезрителей, бесед с родителями я пришел к убеждению, что еще немало 
людей, которых слово «режим» пугает. Правда, с понятием  
«режим» у них чаще всего связывается только регулярный прием 
пищи и соблюдение времени, когда надо ребенка укладывать 
спать. А что касается организации свободного времени  
ребенка, то тут само понятие «режим» они считают неприемлемым, 
даже вредным. 
«Дайте ребенку свободу. Не сковывайте его режимными 
требованиями»». 
«У нас в семье не было никакого режима и слава богу. Все 
мои братья и сестры стали людьми интересными и полезными, 
и все потому, что у каждого было свое дело и он занимался чем 
хотел и когда хотел, и никто нас не ограничивал». 
Чаще всего возражения против режима сводятся именно к 
таким эмоциональным всплескам, призывам к свободе и  
опасениям, что пресловутый режим воспитает ограниченных  
педантов и скучных людей. 
Но, дорогие'родители, а задумывались ли вы над тем, 
причиной скольких семейных неурядиц, конфликтов между 
взрослыми и детьми стало отсутствие здоровой, разумной  
организации жизни в семье, в том числе и свободного времени 
детей? 
Ведь ратуя за режим, мы имеем в виду именно это:  
наилучшим образом помочь физическому и духовному развитию  
ребенка, формированию его индивидуальности. 
Сама природа задает нам ритм жизни. Разве не диктует 
определенный режим смена дня и ночи, утомляемость челове- 
44 
ческого организма, периодическая потребность в еде, сне и 
т.д.? 
И вряд ли кто-нибудь будет оспаривать благотворность  
разумного чередования труда и отдыха, физического и  
умственного напряжения. 
Я думаю, что степень размеренной целесообразности жизни 
человека может служить своеобразным мерилом его культуры. 
Ребенок должен быть включен в здоровый ритм жизни 
семьи, чтобы у него с детства сложились определенные  
привычки и навыки. Когда ребенок еще мал, то он просто привыкает 
к тому стилю жизни, который уже сложился в семье, он его 
воспринимает как естественную норму поведения. Важен сам 
уклад жизни в семье и как поддерживают его взрослые. 
Если бы не телевизор, мы бы не отважились после  
многочисленных рекомендаций врачей и педагогов еще раз  
останавливать ваше внимание на этой азбучной истине. 
Дело в том, что когда ребенку 3—4 года, он уже начинает 
приобщаться к телевизору. В большинстве случаев экран его 
занимает минут тридцать в день, хотя разумным пределом  
просмотра телепередач для этого возраста психологи считают 
15—20 минут. Многое трехлетним непонятно и неинтересно, к 
тому же они быстро утомляются, но движущееся изображение 
продолжает удерживать у экрана более продолжительное  
время. Однако пока что родители в телевизоре не видят никакой 
помехи. 
Но ребенок растет, расширяется круг его интересов, крепче 
становится его организм, и в 5—7 лет он уже просиживает у 
телевизора до полутора часов в день. Но поскольку ребенок 
еще не учится и все его время — свободное, вас опять-таки не 
очень волнует увлечение ТВ. Он смотрит передачи «между 
делом», т.е. просмотры перемежаются играми, прогулками, 
сном. Порой даже радует, что беспокойный малыш не маячит 
перед глазами. 
Но вот сын пошел в школу. У него появляются обязанности: 
ежедневно несколько часов занимает учеба в школе, дома — 
приготовление уроков. Уроки не всегда получаются и не всегда 
45 
привлекают, а свободного времени, когда можно поиграть, 
погулять или чем-либо заняться по своему усмотрению,  
становится меньше. А тут и телевизор покушается на это же время, 
предлагая то мультфильм, то викторину, то кинолутешествие и 
т.д. Что же делать? 
Мы уже отмечали, что очень часто школьник  
(преимущественно это учащиеся 3—4-х классов) решает вопрос просто: 
уроки можно сделать позже, а кроме того, робкое  
предположение, ««авось завтра меня по этому предмету не спросят>», к 
началу передачи перерастает в уверенность, и телевизор  
включен... И тогда родители начинают беспокоиться — ТВ  
становится помехой в учебе, ТВ нарушает не только сложившийся 
режим дня, но и каждодневно подвергает серьезному  
испытанию столь упорно и кропотливо вырабатываемые у детей  
представления об обязанностях и долге. Согласитесь, не так-то 
легко пожертвовать развлечением, а ведь часто этот поединок 
проходит один на один. Вот тут-то режим и его верные  
спутники — привычки необходимы. Только они могут помочь ребенку 
принять правильные и для него самого наиболее полезные 
решения. 
Режим в детские годы способствует формированию черт 
характера, необходимых взрослому человеку. Об этом не нужно 
забывать и не следует уповать на то, что все придет само собой. 
Уменье интересно провести свое свободное время, пожалуй, 
одна из трудных задач даже для взрослого. И, наоборот, 
неуменье активно, с творческой отдачей его использовать 
рождает скуку, неудовлетворенность, какую-то душевную 
вялость. Очевидно, вам не раз приходилось наблюдать, не 
правда ли, бесцельно слоняющихся по улицам подростков, 
бродящих от кинотеатра к кинотеатру, из парка в парк. Они 
ищут, кто их развлечет, кто заполнит их время. У этих ребят 
понятие отдыха, свободного времени связывается с  
представлением о блаженном ничегонеделании. Зародившись в детстве, 
оно потом часто сопровождает человека всю его жизнь. 
Поэтому так важно с ранних лет воспитать у ребенка  
потребность деятельного, творческого отдыха, заполненного интерес- 
46 
ными, ранее намечаемыми делами (пусть это будет выезд в лес, 
прогулка в соседний парк, посещение музея, просмотр  
кинофильма или телепередачи, чтение всей семьей вслух  
юмористического рассказа и т.д.). 
Дети часто не умеют и не знают, чем занять свое время. 
Разве так уж редко, сделав уроки, дочь спрашивает родителей: 
«А что теперь мне делать»?. И как часто в ответ звучит: «Ну, 
почитай что-нибудь, ну, поиграй во что-нибудь»»... 
Но ведь ребенок-то потому и спрашивает, что он не знает, 
что читать и во что играть. У него хотя и есть книги, но  
потребности читать нет, ему хотелось бы поиграть, но он не знает  
интересных игр, а взрослые отмахнулись от него. И ребенок садится 
к телевизору или отправляется на улицу, где без удержу гоняет 
мяч, утомляется и, придя домой, уже ничем не способен  
заняться. Ваша помощь, дорогие родители, так нужна детям. 
Хотя именно, когда заходит речь о свободном времени ребенка, 
более всего противников режима. Но зачем же понимать режим 
как нечто призванное «заорганизовать»» свободное время!? 
Нам кажется, что роль родителей — не в том, чтобы  
регламентировать по часам и минутам, когда и чем ребенку заняться. 
От родителей здесь требуется совет, тактично сделанное  
предложение, как разумно распорядиться своим досугом. 
Не менее существенно, чтобы школьники сами умели  
выделять время и в первую очередь на важные и обязательные дела 
(в том числе для* работы по дому). 
Как часто бывает, когда подросток бросает все свои дела, 
чтобы посмотреть телепередачу. Представим себе, что мать 
или бабушка вместо того, чтобы приготовить обед или ужин, 
садятся смотреть кино. Очевидно, такого не бывает. 
Итак, снова поединок между чувством долга и желанием 
отдохнуть. Что же делать? Как помочь ребенку справиться с 
подобными коллизиями? Думаем, что начинать следует, как  
всегда, с простого. К примеру, через несколько минут должна быть 
интересная передача, а грязная посуда после ужина не убрана, 
не вымыта. Что делает обычно мама или бабушка? Идут мыть, 
пропуская начало передачи, а дети или даже отец (хотя это тоже 
47 
не правило) спокойненько усаживаются у экрана. Это уже  
сигнал опасности. Надо приучать ребят к взаимопомощи.  
Попробуйте разделить обязанности. Вы готовите ужин, а детям  
поручите вымыть посуду, сделать уборку и т.п., а если они вовремя 
не успевают, помогите им. Вот, начиная с таких мелких дел, и 
будет развиваться чувство взаимопомощи. 
Думаю, мы подошли к выводу, что в режиме дня самые 
главные кирпичики — это время, предусмотренное на учебные 
занятия и подготовку к школе, на выполнение домашних  
хозяйственных дел и свободное время. 
Из чего складывается так называемое свободное время? 
При нормальном развитии ребенка большое место занимает 
его самостоятельная творческая деятельность. Разумеется, 
значительным творческим актом рисуется нам в идеале и учеба 
ребят — трудный и радостный процесс познания. Но хочется 
подчеркнуть, что именно в свободные от школы часы дети  
занимаются своими любимыми делами, в которых уже проявляются 
их индивидуальные склонности и способности. Как мы,  
родители, умеем здесь привлечь себе в помощники телевизор? Ведь 
почти каждая передача может быть использована для развития 
творческих интересов детей, может послужить пищей для их 
самостоятельной деятельности. 
Можно считать передачу удавшейся, когда по окончании ее 
начинается «самодеятельность» ребят, будь это простое  
подражание — игра, обращение к определенной литературе, поход на 
ту или иную выставку, или просто обсуждение в кругу друзей. 
Немалую часть своего свободного времени посвящают дети, 
начиная с первых дней школьной жизни, общественной работе, 
выполнению пионерских или комсомольских поручений. Их  
значение велико и как дополнительный источник знаний. Это 
школа определенных умений и навыков, и, самое главное, 
школа становления активной личности будущего гражданина 
нашей страны, формирования его коммунистического  
самосознания. 
48 
Но порой выполнение общественных обязанностей,  
добровольно взятых на себя, сталкивается с непреодолимым  
желанием развлечься, посмотреть интересный фильм. 
Так, в зимние каникулы ребята, участвовавшие в новогодних 
спектаклях, где они с большой охотой и старанием исполняли 
роли снегурочек, зайчиков и даже дедов-морозов,  
отказывались под любым предлогом от выступлений в те дни и часы, 
когда спектакли совпадали по времени с показом  
многосерийного польского кинофильма «Четыре танкиста и собака».  
Многие юные зрители просили заменить их, а находились и такие, 
которые «заболевали» или выискивали неотложные дела,  
придумывали веские причины, чтобы остаться дома и посмотреть 
фильм, кстати, уже неоднократно показанный по ТВ. 
Если вы действительно хотите воспитать сознательного  
гражданина, волевого ответственного человека, то не следует  
проходить мимо подобных конфликтов. Помогите ребенку занять 
правильную позицию. Пусть он испытает удовлетворение и  
гордость от того, что не пожертвовал общественными интересами 
ради личных. 
Занимаясь самовоспитанием, «вырабатывая» характер,  
подростки пытаются проверить себя, придумывая различного рода 
трудности. В данном случае естественный ход жизни  
предоставляет им такую возможность. Не упускайте ее, используйте 
в качестве материала для упражнения и закрепления  
определенных нравственных принципов. 
Думаем, что не стоит оберегать ребенка от всех жизненных 
конфликтов. Наоборот, ищите повод, который может послужить 
тренингом ребенку в выработке волевых качеств. Телевидение 
ежедневно предоставляет массу таких возможностей. Само 
собой разумеется, здесь нужна мера, чтобы не превратить ТВ в 
«сплошной» инструмент для воспитания силы воли. 
Все виды трудовой и общественной деятельности ребенка 
являются весьма существенной частью формирования  
основных качеств его личности. Выше мы уже говорили о той роли, 
которая принадлежит в этом телевидению. Но, чтобы не  
возникло ложного представления, что мы относим телевидение 
49 
только к «средству развлечения», остановимся еще несколько 
подробнее на значении телевидения в учебе ребят. Дело не 
только в том, что по телевидению передаются специальные 
учебные передачи. Большую роль в развитии интересов  
школьников играют различные познавательные передачи типа ««Клуб 
кинопутешествий»», ««Ребятам о зверятах»», ««Всемирный  
следопыт»». Они, несомненно, помогают ребятам в изучении  
географии, литературы или истории. Но их значение гораздо большее, 
чем просто помочь, к примеру, в усвоении темы  
«Пресмыкающиеся»» или ««Малазийский архипелаг»». Встречи с учеными, 
знакомство с нерешенными научными проблемами, с научным 
поиском, с мужественным героическим трудом ученых,  
путешественников, естествоиспытателей — это вовлечение ребят в 
духовный мир научного поиска, научного подвига. Мир знания 
предстает увлекательным, полным приключений, мужества, 
хотя он облекается в форму хроникальных, документальных 
съемок, лишенных острой завязки, развернутого сюжета. Тем 
не менее в опросах, которые проводились среди ребят, именно 
эти передачи оказывались любимыми. 
Заканчивая разговор о первой «азбучной истине»», мне  
хотелось бы рассказать об организации жизни в семье одного моего 
ленинградского знакомого. Зовут его Игорь Петрович. По роду 
занятий у него и его жены Любови Александровны  
ненормированный рабочий день. Это означает, что они приходят домой не 
в определенное время, часто задерживаются на работе. Дети 
после школы хозяйничают дома сами, а их двое: ученики 3-го и 
8-го классов. 
Но по субботам за ужином собираются все вместе, толкуют 
об учебе, о делах на работе, о событиях в стране, о прочитанной 
книге, увиденном фильме — обо всем, что заинтересовало и 
привлекло внимание. 
Эти беседы проходят непроизвольно, непреднамеренно, 
являясь выражением интереса друг к другу. Конечно, подобные 
разговоры возникают и в другие дни. Но в субботу обязательная 
часть этих бесед — чтение и разбор телепрограммы на неделю. 
50 
Отец надевает очки, берет карандаш и читает программу вслух, 
все вместе комментируют ее. 
В непринужденной беседе вместе решают когда кому что 
смотреть. 
Например, отмечаются передачи для малышей, их две ино 
гда три в неделю (по 30 минут). Они предназначены как раз для 
третьеклассника Толи. 
Отмечаются две, не более, передачи для восьмиклассника 
Сережи, но он не всегда смотрит именно их, так как ему  
интересны научно-познавательные, а иногда и литературные  
программы для взрослых. 
Выбираются передачи для всех членов семьи. Это «Клуб 
кинопутешествий»», ««Следствие ведут знатоки»», а также  
наиболее интересные спортивные соревнования,  
литературно-музыкальные передачи, фильмы. 
Все вместе смотрят программу ««Время»». 
После программы «Время»» (не более трех раз в неделю) 
взрослые смотрят вечерние передачи. В субботу старшему 
сыну тоже разрешается смотреть взрослые передачи. 
Так, в результате всестороннего обсуждения интересов, дел, 
развлечений, рождается своя семейная программа на неделю. 
Малышу предоставляется почетная возможность цветными 
карандашами выписать отмеченные передачи в отдельный  
альбом, украсить программу рисунками. В этом же альбоме ребята 
после просмотра пишут отзывы о передачах, иногда делают 
зарисовки, особенно старается младший. Иногда рисует и отец. 
Записи, которые делают мальчики, лишний раз дают им  
возможность задуматься над увиденным, точнее выразить свои 
мысли и впечатления. 
Отец и мать всегда относятся к составлению своей семейной 
«программы» серьезно, хотя обсуждение проходит весело, с 
шуткой, непринужденно. Родители стараются «не давить» на 
ребят, не гасить их желаний, а убеждать, советовать. Ребята 
при составлении «программки»» не чувствуют, что родители 
пытаются ограничить их телевизионные увлечения. Дети  
слышат разумные доводы, и если родители пытаются на чем-то 
настаивать, то стараются обосновывать свои возражения. 
51 
В процессе выбора передач одновременно решаются и  
другие вопросы: когда читать, играть, заниматься своими делами. 
Прочитав всю программу и получив «разведданные»», родители 
спрашивают: «Ну, ребята, давайте вспомним, у кого какие дела 
в школе?» 
Рассматривая телепрограмму недели, семья, по существу, 
не только определяет кто что и когда будет смотреть, а скорее 
разумно прикидывает кто чем будет заниматься, выделяет 
время для покупки продуктов, уборки комнат, прогулок на  
воздухе, общественных дел, спорта, чтения и т. д. Учитываются и 
индивидуальные склонности ребят. Сережа увлекается  
техникой, Толя любит лепить. Это, конечно, принимается во внимание 
при выборе передач. 
Родители отнюдь не ставят перед собой цели заполнить до 
конца все свободное время детей. Ребята имеют право и что-то 
менять, и что-то добавлять. Но это не распространяется на  
телепередачи. Поправки тут вносятся только в виде исключения, 
когда происходят особые события, интересные для всех.  
Конечно, чтобы добиться успеха в подобных начинаниях, нужны 
определенные традиции в семейных отношениях. 
Уже много лет в семье Игоря Петровича один раз в неделю 
составляется веселое, забавное расписание под названием 
«Поможем нашей маме», в котором предусматривается участие 
детей в уборке квартиры, чистке обуви, мытье посуды, покупке 
продуктов и т.п. Причем эти обязанности (конечно, в большем 
масштабе) ложатся, естественно, и на отца. Здесь все давно 
привыкли помогать друг другу, знают свои обязанности,  
привыкли многие вещи делать коллективно, весело; одобрение или 
осуждение, выраженные иногда лишь покачиванием головы, 
словом похвалы, шуткой, с малых лет регулировали поведение 
детей. 
Великий русский педагог К. Д. Ушинский писал, что у  
человека портится голова и сердце, и нравственность, если он не 
знает, как использовать свободное время. 
К сожалению, телевидение несет в себе опасность стать 
52 
самым доступным развлечением, которое требует самого  
минимального личного усилия. 
Смешно было бы в данном случае винить телевидение или 
детей. Все дело в неумении взрослых помочь детям  
организовывать свою деятельность, чередовать игры, увлечения и  
развлечения. 
Итак, «делу — время, потехе — час» — это азбучная истина 
№ 1. 
АЗБУЧНАЯ ИСТИНА № 2 
«Ваше собственное поведение — самая решающая вещь». 
Эта истина сформулирована А. С. Макаренко. Она утверждает 
мысль, что жизнь самих родителей должна быть организована 
таким образом, чтобы служить примером для детей. 
Часто взрослые говорят: «Дети любят подражать». Каждый 
из нас наблюдал, как малыши, играя в «куличики»», повторяют 
движения взрослых, как подрастая, они играют в «дочки-  
матери», в магазин, в больницу, в школу, копируя поведение тех 
взрослых, с кем им приходится сталкиваться (родителями,  
врачами, учителями). 
В играх отражается труд человека, взаимоотношения людей 
между собой. Воссоздаются конфликты, с которыми дети уже 
встречались... Одним словом, они изо дня в день накапливают 
опыт вхождения в самостоятельную жизнь. 
В подростковом возрасте дети, читая книги, мечтают быть 
похожими на героев повестей и романов. 
У ребят старшего школьного возраста появляется свой  
идеальный герой, проявляющий лучшие человеческие качества в 
самых сложных жизненных ситуациях. 
Так что дело не в том, что дети «любят подражать». Это  
естественная психическая деятельность человека, связанная с 
ростом, развитием, овладением социальным опытом людей. А 
начинается она с подражания отцу и матери, потому, что они 
самые близкие люди, и детям кажется, что их мама — самая 
лучшая, а папа — самый сильный и умный. 
Ребята не просто внешне копируют родителей, они разде- 
53 
ляют их мнения, оценки и смотрят на жизнь до определенного 
возраста глазами своих родителей. Взрослея, они, конечно, 
критичнее оценивают родителей, а их идеалом становятся 
более совершенные герои литературных произведений. 
Мо сейчас мы заговорили о роли примера взрослых только 
для того, чтобы подсказать вам, что, если содержание вашего 
свободного времени ограничивается телевизором, если вы 
забыли дорогу в кинотеатр, театр или музей, если вас не  
увлекают лыжи, походы в лес, рыбалка, если вы давно не держали 
в руках книгу, не рисуете и не мастерите, то знайте: ваше 
пассивное времяпрепровождение будет казаться ребенку  
нормальным. И не сетуйте на то, что он не умеет ничем заняться, что 
он как маленький «старичок»» будет рядом с вами тихо и мирно 
коротать время у экрана ТВ. 
И не сетуйте тогда на своего ребенка, если он будет слаб 
физически, не полюбит книгу, будет мало знать. Не думайте, что 
чем больше просмотрит он передач, тем станет умнее,  
образованнее. К сожалению, нет. 
Если вы не взяли за привычку выбирать из многообразия 
программы передачи, которые следует посмотреть, если вы 
считаете, что совсем ни к чему заранее определять, что вам 
необходимо сделать по дому, или когда сходить с семьей в 
гости, а когда — за грибами, когда всем вместе почитать, то 
избирательного, разумного и экономного подхода к  
использованию свободного времени не ждите и от ваших детей. 
Это истина. Не забывайте о ней, и, особенно, когда в семье 
появился телевизор. Роль личного примера родителей в их 
отношении к ТВ имеет решающее значение. 
АЗБУЧНАЯ ИСТИНА № 3 
«Каждая передача должна стать событием». 
Да, и для тех, кто ее создает, и для тех, кому она  
предназначена. Сомнений нет относительно тех, кто ее создает,—  
телевизионщиков. Прежде чем выдать передачу в эфир, надо ее за- 
54 
думать, спланировать тысячу вещей для ее реализации. В  
подготовке передачи участвуют десятки человек. 
Немало приходится потрудиться инженерно-техническим, 
и творческим работникам. 
К сожалению, для иных зрителей не каждая передача  
является событием, а хотелось бы, чтобы таковой становилась 
каждая, на какой бы возраст она ни была рассчитана: 
дошкольный или школьный. 
Художник В.Курчевский показывает, как нарисовать с  
помощью двух-трех кружочков утку, курицу и индюка. И для ребят 
это — открытие. Это событие. 
Диктор Валентина Леонтьева «по секрету»» рассказывает 
детям, какие можно приготовить к 8 Марта подарки мамам, 
бабушкам, сестричкам. И этот секрет — событие для малышей. 
В одной из передач «Костер»» пионеры узнали о том, как 
интересно, по-настоящему интересно идет жизнь отряда одной 
из ленинградских школ. 
Юные историки и техники, поэты и любители музыки,  
посмотрев предназначенную для них передачу, что-то впервые  
узнают, открывают для себя. 
Старшеклассники участвовали в математической  
олимпиаде, они решили много оригинальных задач, перед ними  
выступали крупные ученые, и это событие в их жизни. Все может, 
должно быть событием, но всегда ли становится им? 
Не будем касаться качества передач, это не входит в нашу 
задачу. Но, помимо качества, есть условия, которые надо 
соблюдать юным телезрителям, чтобы каждая передача  
оставила в их душе след. 
Родители ни в коем случае не должны спокойно относиться 
к тому, что дети смотрят передачу, не рассчитанную на их 
возраст, уповая, что передача вреда не принесет.  
Действительно, прямого вреда не принесет. Но не принесет и пользы. Более 
того, ребенок потеряет время впустую, снизится уровень его 
восприятия. Передача не вызовет необходимого душевного 
отклика. 
Следующее необходимое условие — это учет  
индивидуальных интересов детей. И здесь закономерно встает вопрос: все 
55 
ли передачи, рассчитанные на возраст вашего ребенка, следует 
ему смотреть? 
Конечно, есть передачи общественно-политического  
характера, которые интересуют многих ребят, независимо от их  
личных склонностей. Это передачи на школьно-пионерские темы, 
возможно, это программа «Подвиг», «Время» или музыкально- 
развлекательная и познавательная игровая передача  
«Будильник». 
И все же — если передача по своей тематике  
предназначается для юных историков или техников, любителей спорта, 
музыки и природы, то она станет событием в первую очередь 
для тех, кто ее ожидает: историков, техников, спортсменов и 
природоведов. 
Одним из наиболее важных условий эффективности  
передачи мы считаем привычку выбирать передачу, заранее  
внутренне подготовившись к ней, освободив для нее время. Нет 
ничего хуже, когда ребята смотрят три-четыре передачи подряд, 
не делая перерыва, не осмысливая увиденное, не  
задерживаясь на нем. 
Помню один неудачный случай верстки программы. Шла 
передача о Брестской крепооти. Зрители стали свидетелями, 
как население Бреста встречалось с защитниками крепости, 
как под звуки траурной музыки возлагались цветы на могилы 
погибших. Передача кончилась, вызвав слезы на глазах и у 
зрителей, и у тех, кто ее готовил. А через 30—40 секунд из другой 
студии другая бригада (но по тому же каналу) стала выдавать 
очень веселую передачу для молодежи «Учитесь танцевать». 
Это, конечно, была грубая ошибка программирования, явно без 
учета того, что еще многие зрители смотрят передачи подряд, 
без перерыва. Одно впечатление накладывается на другое,  
нейтрализует его, и зритель отходит от телевизора, не унося 
ничего, кроме головной боли да зрительного утомления. 
А теперь для сравнения вспомните, какое огромное  
удовольствие вы получаете от просмотра хорошего фильма или 
спектакля. Но если кому-то из вас довелось присутствовать на 
кинофестивалях, длящихся по две-три недели, где ежедневно 
56 
приходится смотреть минимум два фильма, вы помните, в какую 
пытку превращается это обилие впечатлений. 
Подобное часто происходит перед телевизором. Школьник 
смотрит передачу за передачей. Перед его глазами проходит 
целый поток впечатлений, идей, представлений. Еще не  
осмыслив одно, он уже включается в переживание другого, третьего. 
Что же остается в его памяти? В его сердце? 
Такое отношение к кино А. С. Макаренко называл ««пустым 
глазением». Как бы ни была хороша каждая передача в  
отдельности, при подобном просмотре она не задевает ребенка, не 
оказывает на него благотворного влияния, делает его  
пресыщенным, равнодушным существом. 
Родителям очень важно с детства воспитывать у ребенка 
зрительскую культуру. И взрослым надлежит тут быть первым 
примером. 
Наконец, последнее условие: родители сами должны  
стремиться к тому, чтобы сделать из передачи событие для ребенка, 
пусть маленькое, но событие. Мы уже говорили, как для этого 
важно заранее наметить передачу, подготовиться к ней, 
создать у юного зрителя настроение нетерпеливого ожидания, 
освободить специально для этого время, а во время просмотра 
передачи — не мешать ему — не отрывать посторонними  
вопросами, не поручать каких-либо дел и т. д. 
Если у вас есть возможность, присядьте рядом с ребенком и 
вместе с ним посмотрите передачу. Вы даже не представляете, 
как это нужно. И не только ему. Но и вам. Совместный просмотр, 
сопереживание, обмен мнениями, имеют большое  
воспитательное значение. Вы переживете счастливые минуты, когда  
вечером или на другой день ваши дети, перебивая друг друга, начнут 
рассказывать кому-то третьему об этой передаче и вы узнаете 
свои мысли, и как преломилось в сознании ребенка то, на что 
вы обратили его внимание, как он учится выражать свои мнения 
и чувства. И как накапливается его жизненный и  
художественный опыт. 
Но вот что рассказывают некоторые родители: ««А мой сын 
сердится, если я начинаю его допрашивать после переда- 
57 
чи».Очень верное наблюдение. Что же делать? Во-первых, не 
надо «допрашивать», а во-вторых, впечатления ребят бывают 
настолько сильны, что они часто чувствуют себя буквально 
подавленными всем увиденным. Им не сразу удается  
разобраться в самих себе. Подождите, не спешите. Завтра или, может 
быть, даже через несколько дней вы обязательно узнаете, что 
так взволновало и поразило вашего мальчика или девочку. 
Только будьте терпеливы и тактичны. И как бы ершисто ни 
приняли они ваше предложение высказать свое мнение, 
можете быть уверены, они не пропустят ни слова из ваших 
суждений. Заведите после передачи разговор между  
взрослыми, как вы это делаете после просмотра кинофильмов или 
театральных спектаклей, интересных экскурсий, походов в лес. 
Обсуждайте, высказывайте свои оценки, включайте в обмен 
мнениями «на равных» и своих детей. Незаметно у них будет 
вырабатываться свой вкус, свои критерии. 
Социолог Н. С. Мансуров очень удачно сказал, что семья, ее 
мнение — это своего рода «точка отсчета», на которую  
ориентируется ребенок, встречаясь с другими людьми, с различными 
жизненными ситуациями. 
Сегодня мы разговариваем- с родителями, которые хотят 
сохранить любовь и доверие своих детей на всю жизнь. 
Мне хотелось бы в связи с этим напомнить, что Карл Маркс 
этого добился потому, что он не исключал из сферы своей 
духовной жизни дочерей даже тогда, когда они были еще  
крошками. Они вместе с отцом и матерью переживали трагедию 
коренного населения Америки, зверскую эксплуатацию негров 
и т. д. Одна из дочерей писала письма президенту США 
Линкольну, «советуя», как тому надо поступать, а К. Маркс  
регулярно «носил» их на почту. Конечно, дочери не могли  
разбираться в сложной политической борьбе, но они хорошо  
чувствовали эмоциональную направленность, царившую в семье. 
Карл Маркс, говоря о рабочей печати, подчеркивал, что она 
приносит в каждую хижину «дух государства». Это определение 
целиком можно отнести и к телевидению. «Дух государства», 
находясь в руках родителей, может значительно повысить 
содержание воспитательной деятельности семьи. 
58 
АЗБУЧНАЯ ИСТИНА Nt 4 
«Чтение — вот лучшее учение» 
Хрестоматийными стали эти слова А. С. Пушкина. 
Можно ли представить себе передачу эстафеты знаний,  
опыта, культуры от поколения к поколению без книги? Но вот 
появился в нашей жизни, в нашем доме телевизор, и монополия 
книги как основного источника знаний явно поколебалась. 
Сказку можно теперь не только прочитать, но и услышать в 
исполнении хорошего актера или увидеть в мультфильме. Даже 
с учебниками телевидение пытается вступить в соперничество. 
В школах все чаще практикуются телеуроки. В некоторых  
капиталистических странах нехватка школ и учителей даже привела 
к созданию телешкол. 
Широкое развитие телевидения посеяло кое у кого тревогу: 
не умрет ли печатная книга после 500-летнего своего  
существования? Более того, появились теории и теоретики,  
предсказывающие ее неминуемую гибель в век электроники. Под влияние 
подобных теорий подпал один из прогрессивных  
западногерманских поэтов Энценсбергер, который утверждает, что с  
развитием социального прогресса получат признание более  
демократические средства общения, такие, как радио и телевидение, 
устное поэтическое творчество. Книги останутся для узкого 
круга читателей или как служебная литература для отдельных 
специалистов. 
Еще радикальнее высказывается о смерти книги канадский 
профессор Маршалл Макклюэн, объявивший «закат галактики 
Гутенберга» и наступление новой электронной эры, эры  
возрождения и господства устной речи и т. д. 
В нашем веке книге уже дважды приходилось выдерживать 
подобные атаки: со стороны радио и кинематографа. Как вы 
знаете, с изобретением кинематографа и театру пророчили 
неминуемую гибель. Однако кино не погубило театра. Театр же 
использовал некоторые кинематографические приемы и  
принципы, расширив свои художественные возможности. С  
появлением телевидения гибель и театра, и кинематографа казалась 
лишь вопросом времени. Но прошли годы, десятиле- 
59 
тия — ««живы оба», более того, развитие телевидения даже  
способствовало тому, что театр стал более ««театральным», а  
кинематограф — ««кинематографичным». 
Что же, может быть, и книга обрела нечто новое?.. 
Еще Стефан Цвейг, выступая против предсказаний о близкой 
гибели книги в век технического прогресса, писал: ««Ни одному 
источнику энергии не удалось создать такого света, который 
исходит порой от маленького томика, и никогда электрический 
ток не будет обладать такой силой, которой обладает  
электричество, заложенное в печатном слове. Нестареющая и  
несокрушимая, неподвластная времени, самая концентрированная 
сила в самой насыщенной и многообразной форме — вот что 
такое книга, так ей ли бояться техники, разве не с помощью тех 
же книг техника совершенствуется и распространяется? 
Повсюду, не только в нашей личной жизни книга есть альфа и 
омега всякого знания, начало начал каждой науки». 
Границы размежевания между книгой и другими средствами 
массовой коммуникации осознаются в наши дни довольно 
определенно: ««Радио и телевидение можно назвать средствами 
«моментальной» или ««экспресс-информации». Их  
информационное действие прекращается одновременно с  
прекращением радио или телепередачи. Книги же — средство  
постоянной информации, это огромный резервуар, из которого можно 
черпать информацию в любой момент и в любом  
количестве»,— пишет Л. Владимиров. 
Бурное развитие телевидения в наши дни сопровождается 
одновременно не менее стремительным подъемом книжного 
дела. Казалось бы, что телевидение и радио, поднимая  
культурный уровень народа, уже тем самым способствует развитию 
потребности к чтению. 
Однако каждая семья, получая информацию по нескольким 
каналам, порой пользуется ими стихийно. Часто и взрослые 
заполняют все свое свободное время просмотром телепередач, 
не говоря о детях, которые еще скорее, чем взрослые, готовы 
пожертвовать чтением, радиопередачей. 
60 
Не буду повторять, как важно, получив в распоряжение 
источники познания и культуры, воспитать у детей  
избирательное отношение к ним, так, чтобы они дополняли друг друга. 
Здесь надо обратить ваше внимание, что детское  
телевидение постоянно стремится укрепить у детей привычку к чтению, 
знакомить с новыми книгами, писателями, рассказывать об 
истории создания той или иной книги. 
Но повторяю, в первую очередь от вас, товарищи родители, 
зависит, как разрешит ваш ребенок вопрос о соперничестве 
между книгой и ТВ. 
Ученые заметили, что в «читающих» семьях дети охотнее 
читают книги. 
И если мы хотим, чтобы наши дети любили читать,  
существует совершенно простое правило, доступное любому  
родителю: читайте детям книги, и даже тогда, когда ваш малыш 
пойдет в школу и сам начнет читать. Старайтесь, чтобы ваш 
ребенок пересказывал прочитанное. 
Не все дети одинаково быстро и хорошо научаются читать. И 
вот для тех, кому чтение дается с трудом, сам процесс чтения, 
являясь в школе обязательным, становится неприятным. Уже с 
1-го, 2-го классов у некоторых детей пропадает интерес к 
чтению. Это очень опасный сигнал, который родители должны 
предупредить. Детское любопытство, интерес к содержанию 
сказки, рассказа торопит ребенка вперед, а он движется очень 
медленно, много усилий уходит на овладение техникой чтения. 
Тут-то телевидение выступает «счастливым соперником». 
Мне пришлось встретиться с семьей, в которой ребята 4-го 
и 6-го классов совершенно равнодушно относились к книгам. 
Родители надеялись: подрастет — сам начнет читать. Но вот 
миновал четвертый, пятый, шестой класс, внимание детей  
целиком захвачено телевидением. Кроме того, у ребят явно начали 
проглядывать интересы к технике. И тогда родители решили, 
что больше ждать нельзя. Они постепенно установили традицию 
читать вслух интересные сообщения из газет, журналов. Если 
кто-то из взрослых, читая какую-либо книгу, находил  
интересное место, он зачитывал его для всех и все вместе обсуждали 
61 
прочитанное. И так, естественно, члены семьи включались в 
чтение. А начиналось обычно с предложения: «Вчера в  
«Литературной газете» встретился очень смешной рассказ в разделе 
«Клуб 12 стульев»». Хотите прочту?»» За отдельными рассказами 
и фельетонами последовало чтение отрывков из романа  
Ярослава Гашека «Похождения бравого солдата Швейка», потом 
засели за книги И. Ильфа и Е. Петрова. Особенно с большим 
успехом прошло чтение «с продолжением» повести ф.  
Искандера «Созвездие Козлотура». Привлекла внимание семьи и  
драматическая повесть «А зори здесь тихие»» Б. Васильева, и даже 
«Обломов» И. А. Гончарова был прочитан вслух, а потом «Пикк- 
викский клуб» Диккенса, произведения Станислава Лема. 
Чтение было не ежедневным, но систематическим, чаще 
всего по субботам и воскресеньям, летом во время отпуска 
читали почти каждый день. . 
Читали не только для детей, их участие в слушаньи было 
добровольным. Родители приносили домой для себя (или как бы 
только для себя) то научную фантастику, то захватывающий 
детектив, то повести и романы. Читали их «по кругу» (причем 
родителям, естественно, многое приходилось читать не в  
первый раз). В результате через 1,5—2 года у ребят появился  
самостоятельный интерес к чтению. Теперь уже они сами начали 
приносить домой книги. 
А самое главное, любовь к чтению стала для детей заметным 
регулятором в их отношениях с телевизором. Теперь они не 
усаживались бездумно к экрану. Они выбирали: «Это надо 
посмотреть. А на это не хочется терять время — у меня лежит 
интересная книжка». 
Ребятам, знающим цену книге, не угрожает телемания. Но не 
слишком ли много слов в защиту книги? Нет, немного. ТВ 
заметно осложнило ей жизнь. У нее сейчас трудный период. Ей 
надо помочь. А она в свою очередь сделает ваших детей 
умными, образованными и воспитанными людьми. 
62 
АЗБУЧНАЯ ИСТИНА Ni 5, 
или «Охота пуще неволи» 
В вагон трамвая входит мама с пятилетним мальчиком. Юный 
пассажир сразу же устремляется к окну. Это не каприз, это 
потребность смотреть, видеть, узнавать. Окно становится как 
бы экраном, через который он видит улицу, людей. Причем, 
ограниченное рамой, окно как бы уменьшает поле зрения, зато 
сосредоточивает внимание на отдельных предметах. 
— Ма! Смотри, смотри. Ма, собака. Ма, смотри, машина. 
Почему она стоит? 
Мама смотрит и видит и собаку, и машину... и ничему не  
удивляется. 
Ребенку все интересно, он еще очень мало знает. Интерес 
руководит им и тогда, когда он пробует на вкус незнакомую 
ягоду или ломает игрушку и когда на обоях пытается изобразить 
чернилами человека или самостоятельно зажечь газовую плиту, 
вбить в пианино гвоздь. 
Достается юным исследователям и изобретателям за их  
опыты! 
Родители пытаются унять не «в меру» любопытных детей, 
успокоить с помощью развлечений, — занять какой- либо 
деятельностью. Усилия родителей не пропадают даром.  
Познавательные интересы детей или приглушаются или...  
развиваются, укрепляются, дифференцируются. Мы даже склонны порой 
несколько преждевременно говорить: «Он будет строителем, 
автомобилистом, историком, врачом, художником...» 
Но бывает и по-другому. Родители не знают ничего об  
увлечениях своих детей. Да и дети, оканчивая десятый класс, так и 
не знают, что же их интересует, куда идти дальше. 
Интерес играет огромную роль в развитии личности  
человека, его деятельности, формировании его характера. Начинает 
интерес проявляться у человека в самом раннем возрасте как 
естественная потребность, как выражение отношения человека 
к окружающему миру. 
Вначале ребенок выражает интерес к предметам, к людям 
потому, что они ему «нужны», эмоционально привлекательны. 
63 
Интерес становится мотивом, побуждающим детей к действию, 
к самостоятельной деятельности (игра тоже деятельность) и 
активному поведению. Эта деятельность добровольна и  
направлена к удовлетворению детского интереса. В процессе  
деятельности развиваются интеллектуальные, нравственные,  
эмоциональные и волевые качества маленького человека. И здесь, мы 
снова это подчеркнем, имеет огромное значение семья. 
Какие игрушки вы покупаете ребенку? Как вы пытаетесь 
успокоить, занять его, чтобы он не мешал вам? Наказываете вы 
ребенка или нет, если он захотел посмотреть устройство 
будильника и разломал его? Все это имеет прямое отношение 
к развитию интересов детей. А теперь ответьте еще на один 
вопрос. Как вы используете телевизор для развития интересов 
ребенка? От правильного ответа зависит, как в дальнейшем 
«поведет себя» ТВ. Он может и «глушить»» интересы детей и 
развивать их. 
В одной семье ребенок рисует, конструирует, играет в  
футбол, ухаживает за цветами или рыбками. 
В другой семье — ребенок того же возраста. Он столь же 
любознателен, но его интересы удовлетворяются... с помощью 
телевизора. 
Он «приобщается» к футболу и другим мужественным видам 
спорта через экран. Он знает о цветах может быть даже и 
больше, чем тот мальчик, который выращивает их сам. Он по 
телевидению знакомится и с юными конструкторами, и юными 
художниками. Но он сам ничего не лепит, не строит машин, не 
поливает и не пересаживает цветов... 
Казалось бы, что жажда познания, столь естественная у 
детей, удовлетворяется в обоих случаях. Только в первом  
случае познание жизни идет деятельным, активным путем. Во  
втором — у ребенка создается иллюзия сопричастности и к  
спорту, и к миру знаний. 
В первом случае ребенку в процессе игры, занятий  
приходится проявлять (а значит упражнять, закреплять) различные 
нравственные качества: настойчивость в достижении цели,  
коллективизм в играх, справедливость и доброту в отношении к 
64 
товарищам, не говоря уже о выработке необходимых  
профессиональных навыков, сноровки, умений. 
Во втором случае ребенок пассивен, свои знания он не 
упражняет в практической деятельности. Он и себя плохо знает, 
ибо не рисует, не моделирует, не пробует своих сил и  
возможностей. 
Итак, свойства личности развиваются в процессе жизни и 
деятельности человека. 
Окружающая социальная среда только в том случае  
обладает воспитательной силой, если она затрагивает потребности, 
интересы ребенка и если у него складываются определенные 
деятельные, практические отношения с ней. Чем  
разностороннее связь у подрастающего человека с окружающим миром, 
чем шире вовлекается он в практическую деятельность, тем 
глубже познает он мир, тем успешнее раскрываются его  
духовные силы, а сам он быстрее формируется как личность. 
Нельзя научиться кататься на коньках, работать на станке, 
создавать произведения искусства по телевидению, без  
обучения на практике, без соприкосновения с жизнью. 
Нельзя развить у детей и нравственные качества,  
предоставляя им лишь возможность смотреть спектакли, телефильмы на 
морально-этические темы. Этого мало. 
Социологические исследования, проведенные кандидатом 
философских наук Э. С. Соколовой, подтверждают, что в тех 
семьях, где досуг школьников был заполнен преимущественно 
пассивными развлечениями, неумеренными телепросмотрами и 
только изредка практическими играми и занятиями, там ребята 
меньше интересовались трудовой и общественной  
деятельностью. 
Мы хотим, чтобы наши дети выросли людьми деятельными, а 
не пассивными созерцателями, взирающими на мир только 
через стекло экрана телевизора. 
Но как быть, если у вашего сына или дочери нет никакой 
особой страсти, никакого увлечения? Они не мастерят, не  
рукодельничают, не посещают кружков в школе или Домах  
пионеров, не играют со сверстниками ни дома, ни на улице. Чаще 
I 
65 
всего это как раз те ребята, которые и попадают в крепкие 
объятия телевизора. 
Семья должна развивать, расширять и направлять интересы 
детей. Чтобы развивать интересы, нужно создать дома  
определенные условия: отвести «рабочее место» ребенку, помочь ему 
достать нужные материалы, оборудование. Этому послужит и 
новая игрушка, и прогулка родителей с детьми, организация 
домашнего уголка-мастерской и включение детей в  
разнообразную и особенно, в коллективную деятельность. Интересы 
проявляются у ребенка очень рано, и большую роль здесь, как 
и во всем, играет подражание и, естественно, тем людям, кто 
ему ближе всего. И вы сами не раз это наблюдали: если отец 
любит слесарничать — и сын тянется за ним, мама шьет — и 
девочка обшивает своих кукол. 
Детям частенько достается от родителей: «Что же ты  
забросил своих бабочек?»... «Сил нет никаких: вчера какие-то жуки, 
сегодня черепки, лишь бы мусор тащить в дом». Не журите 
своих ребятишек, когда они переходят от одного увлечения к 
другому; это естественно в их возрасте. Пусть они попробуют 
себя в самых разных областях. 
Вряд ли поступают правильно и те родители, которые с 
ранних лет предрешают будущее своего ребенка, пытаются 
сделать из него музыканта или художника, инженера или врача. 
Преодолевая иногда страшное сопротивление, на это уходят 
все силы, все средства. В конечном итоге иногда действительно 
получается средний или даже хороший музыкант или художник, 
но, может быть, погублен талант хирурга или биолога. 
Поэтому, воспитывая своих детей, не сужайте их интересы, 
не «специализируйте» их раньше времени, дайте им пошире 
познакомиться с человеческой деятельностью во многих  
отраслях знаний. 
В современной педагогике и психологии проблеме интереса 
уделяется очень большое внимание. Многие ученые считают, 
что именно интерес является стержнем развития личности, тем 
свойством души , которое толкает человека на преодоление 
66 
всевозможных трудностей, делает его одержимым своим 
делом. 
Чем полнее раскрыты интересы детей, чем больше они 
связаны с проявлением детской энергии, фантазии, умений и 
эмоций, тем более осмысленно и разборчиво дети подходят к 
выбору телепрограммы, тем больше дает им просмотр передач, 
углубляя и расширяя их интересы и знания. 
Советские ученые, педагоги, социологи вывели  
закономерность, что дети, занятые в школе общественными делами, 
активно участвующие в работе кружков, в деятельности  
пионерской организации или комсомола, в спортивных секциях, 
проявляют большую самостоятельность в выборе  
телевизионных программ, чем те ребята, у которых нет устойчивых  
интересов или общественных обязанностей и поручений. 
«Общественники-» решительнее отказываются от просмотра 
спортивных соревнований ««не главных команд», от передач, не 
представляющих для них большого интереса, так же как и от 
других занятий, внешне заманчивых, но отвлекающих их от 
учебных занятий или любимых дел. < 
Подобное избирательное отношение — показатель  
целеустремленности личности, ее готовности к самовоспитанию. 
Итак, истина пятая: любимое увлечение юного натуралиста 
или техника, юного историка или спортсмена, редактора  
стенгазеты или коллекционера, пионервожатого или юного артиста 
действительно сокращает пустое времяпрепровождение,  
способствует избирательному отношению к телепрограммам и 
более целенаправленному использованию свободного  
времени. 
Но не только семья, детский сад или школа способствуют 
развитию интересов детей. Само телевидение уделяет этому 
много внимания. Оно в состоянии не только возбуждать в детях 
интерес, но и поддерживать его. Отдельные 20—30-минутные 
передачи могут стимулировать самостоятельную деятельность 
детей, которая потом продолжается дни, недели, а иногда  
становится «постоянно действующим»» независимо от ТВ  
увлечением. В первую очередь это передачи-конкурсы, передачи-викто- 
67 
рины, которые побуждают детей к поискам ответов на вопрос, 
и т. д. Но о них речь пойдет в следующей главе. 
В заключение, уже как бы по традиции, мы попытаемся 
задать себе несколько вопросов. 
Что мы делаем, чтобы широко раскрыть мир перед глазами 
детей? 
Как мы включаем детей в этот мир? Вышли ли вы вместе со 
своими жильцами и детьми на субботник по уборке двора и 
посадке деревьев? 
Помогали ли вы собирать сыну или дочери макулатуру или 
металлолом или, наоборот, делали это нехотя, с иронией  
взрослого над детской забавой, лишь отдавая дань пионерской  
традиции: ««На, мол, возьми, чтобы место не занимало, а то некогда 
было самим выбросить на свалку»... 
Пытались ли вы вместе строить, мастерить, вести календарь 
природы, живо интересуясь тем, что и как делает ваш ребенок? 
Пытались ли вы помочь дочери или сыну подходить к выбору 
программы телевидения с точки зрения развития и углубления 
их интересов, побуждения детей к «обратной связи»,  
практической реакции на передачу? Посоветовали дополнительно что- 
либо почитать по теме передачи, сходить в музей, зарисовать, 
записать в дневник, собрать какой-либо материал, сделать 
модель? 
А если передача спортивная, то помогли ли вы претворить в 
жизнь советы мастеров спорта, если они давались в передаче? 
Или вы прошли мимо всего этого? 
В состоянии ли вы увлечь детей, помочь им найти интересное 
дело? 
ГЛ АВ А 
КОТОРАЯ КАК КОМПАС, ДОЛЖНА ПОМОЧЬ 
ЗРИТЕЛЮ ОРИЕНТИРОВАТЬСЯ ПО КАРТЕ, 
НАЗЫВАЕМОЙ ТЕЛЕПРОГРАММОЙ. 
Очевидно, каждая семья, впервые купившая 
телевизор и не познавшая тайн телепрограмм, чувствует себя 
как группа туристов, заблудившихся ночью и не умеющих 
пользоваться компасом, картой и звездным небом. Где восток? 
Где север? В каком направлении идти? И где же та Полярная 
звезда, по которой можно определить все стороны свОта? 
Конечно, я несколько преувеличиваю, ибо публикуемая в 
газетах программа телевидения отнюдь не звездное небо.  
Многие ориентиры весьма определенны: «Новости», передачи для 
детей, ««Творчество юных»', «Сельская страда», «Программа 
телефильмов»», «Спокойной ночи, малыши»» и т. д. И все же, для 
неискушенного зрителя программа содержит много вопросов. 
Например, что это за передача «Творчество юных»»? Можно 
предположить, что в этой передаче будет рассказано о  
творчестве детей, но о каком именно творчестве? Юных музыкантов 
или авиамоделистов? Юных художников или поэтов? А в каком 
жанре? Концерт юных певцов или киноочерк о юных генетиках, 
получивших новый сорт злаковой культуры? Смотр  
художественной самодеятельности или конкурс рисунков на  
асфальте? 
Мало что говорит и рубрика «Программа телефильмов»». 
Каких фильмов? Документально-производственных, видовых, 
игровых и т. д.? Новых фильмов или повторных, советских или 
зарубежных? Все это остается загадкой. Даже в рубрике  
«Спокойной ночи, малыши» может быть сказка-мультфильм или  
разговор тети Вали со своими кукольными друзьями. 
69 
Каждая программа может оказаться ребусом. В один из 
дней зритель мог посмотреть передачу «А ну-ка, девушки»,  
«Театр одного актера», фильм «Большой и маленький Клаус». Эти 
три передачи шли почти в одно и то же время 21-00, 21-30 и 21- 
45 (не «стрелять» же по программам?). А может быть, вообще 
сходить в этот вечер в кино? 
Но вам кто-то нашептывает, наверное, этакий  
фантастический Телевичок: «В кино можно и завтра, а по ТВ вдруг будут 
инсценированные рассказы Гоголя или Чехова или...» И вы 
остаетесь, смотрите, переключаете с программы на программу 
и в конце концов ругаете себя, что зря отказались от похода в 
кинотеатр. 
Какой же выход из этих блужданий с программы на 
программу попыток проникнуть в смысл передач по названию, 
ожидании чего-то интересного? Каждому телезрителю хочется 
провести вечер так, чтобы не было жаль потерянного зря 
времени. Вот тогда-то зритель начинает сам постигать  
структуру программ, принципы планирования, характер рубрик. 
К тому же ориентироваться ему помогают дикторы,  
объявляющие утром программу на день и по окончании вечерних 
программ на следующий день. Наконец, печатные бюллетени 
заранее сообщают о начале новых циклов, их ведущих и  
авторах, о наиболее актуальных и интересных передачах, а также о 
тех передачах, название которых не определяет их конкретное 
содержание. 
Что же из себя представляет телепрограмма? 
Это, во-первых, не случайный перечень, набор  
кинофильмов и телеспектаклей, спорт-соревнований и новостей. 
Она совершенно не похожа на репертуарный план театров 
или кинотеатров, которыми люди пользуются эпизодически, 
время от времени. 
Телевидение рассчитано на каждодневную встречу, оно 
имеет своих зрителей в каждом доме, в каждой семье. 
Это постоянство аудитории, постоянство встреч  
предполагает необходимость определенной системы в составлении 
70 
передач, связывание отдельных передач в серии, циклы. Циклы 
имеют определенную периодичность. 
Экспериментальным путем установлено, что «пределом 
вместимости» потока передач с экрана телевидения для  
взрослого человека является полтора-два часа (30 минут научно- 
популярные, публицистические передачи и 60—90 минут  
художественные, развлекательные). 
Как строится план передач на день? Так же, как уроки в 
школе. Вначале потруднее, требующие большего  
мыслительного напряжения — математика, физика, химия и т. д. и к концу 
дня — 4—5—6-ые уроки физкультура, пение, труд и т. п. 
Кроме того, учитываются и те часы, в которые больше 
смотрят передачи те или иные категории зрителей, в  
зависимости от рода занятий (школьники, рабочие, студенты, домашние 
хозяйки, воины, колхозники, служащие, пенсионер). 
Поэтому в утренние часы ставятся «новости»», утренняя  
гимнастика, медицинские советы, передачи по домоводству и  
кулинарии, а также повторные передачи для тех, кто не мог их 
видеть в вечерние часы из-за работы или учебы во вторую 
смену. 
Вечером, когда все собираются дома, аудитория  
телевидения становится массовой, то и передачи для взрослых менее 
дифференцированы, их практически могут смотреть все группы 
зрителей. 
Но, поскольку все равно невозможно смотреть 5—6 часов 
подряд, ТВ стремится сделать передачи тематически  
разнообразными, предлагая наряду с художественными,  
познавательными и короткие публицистические очерки, представляющие 
интерес для зрителей различных возрастов и категорий. 
Детские передачи Центрального телевидения имеют 
довольно постоянное место в программе. Правда, если  
посмотреть программы за последние 10 лет, то можно заметить, как 
иногда происходили «сдвижки»» времени детских передач, но 
они, как правило, были незначительными. 
Последние годы объем детских передач, которые готовит 
Центральное телевидение по первой программе, равен 25 
71 
часам в месяц, это с учетом различных возрастных категорий. 
Что же представляет собой детское телевидение, каково его 
содержание? Прежде всего — это система передач,  
рассчитанных на детей дошкольного, младшего, среднего и старшего 
школьного возраста. 
Итак, давайте подробнее рассмотрим программы  
Центрального телевидения, которые выходят в эфир под общей  
музыкально-графической заставкой «Телестудия «Орленок». 
Значительное место в детском вещании занимают передачи 
для малышей, или, как их иногда называют, передачи «Для 
самых маленьких». К этой категории зрителей мы относим 
детей 4—6 лет, хотя приходится иногда наблюдать, как  
трехлетки пытаются, взяв свой стул, пристроиться в первых рядах 
зрителей у домашнего экрана. Однако иминтереснее разбирать 
картинки, неподвижные изображения с комментариями  
взрослых и поэтому через 5—10 минут они находят себе какую-либо 
игру. 
Зато пяти-шестилетние ребята составляют очень  
устойчивую категорию зрителей. Это возраст «почемучек», возраст, 
когда формируется любознательность, наблюдательность,  
пытливость. Примечательно, что желание знать «что-то»  
возрастает, если уже имеются какие-то знания об интересующем  
предмете. У дошкольников почти никогда не возникает вопросов по 
поводу неизвестных предметов. Интерес еще более  
усиливается, когда эти знания эмоционально окрашены. 
В кукольной сказке или инсценированной песенке,  
«Загадке-отгадке» или «Полочке-читалочке», в «Занимательной 
азбуке» или «Ателье «Сарафанчик», в «Гордом кораблике» или 
«Календаре природы» дети узнают о 100 тысячах «почему?»; об 
окружающих в повседневной жизни вещах, о том, кто их делает 
и зачем. А сколько удивительных сведений о природе! Трудно 
найти передачу, в которой бы ребенок не узнавал «Что такое 
хорошо, и что такое плохо...» 
«Выставка Буратино» и «Умелые руки», существующие 
около пятнадцати лет, неизменно вызывают у детей бурную 
деятельность: на студию телевидения нескончаемым потоком 
72 
идУт рисунки и детские поделки. Новые и новые поколения 
втягиваются в полезные игры, в отгадывание загадок, в занятия 
ручным трудом, проявляют интерес к музыке, стихам, чтению. 
Учитывая психологию дошкольного возраста, детские 
редакции избегают обилия незнакомых лиц. Обычно перед  
ними появляются одни и те же дикторы: «тетя» Валя и «тетя»  
Нина, постоянные кукольные герои (раньше — Шустрик и Мям- 
лик, теперь — заяц Тепа, Буратино, песик Филя и мальчик 
Ярошка). 
У ленинградских ребят есть свой любимец — обезьянка. 
Жаконя и другие постоянные герои и друзья малышей. Ребята 
настолько привыкли к ним, что считают их своими близкими 
знакомыми, друзьями. 
Следующая возрастная группа детей — это уже школьники 
7—9 лет, сначала октябрята, а с 9 лет пионеры. Они охотно 
смотрят передачи для малышей, например ««Ребятам о  
зверятах» (трансляция из Ленинграда), «Музыкальную шкатулку»'(из 
Перми), передачи типа ««Играйте с нами», ««Что тебя окружает», 
«Телезнайка». «Театр ««Колокольчик», но одновременно их  
интересует и октябрятская жизнь (««Светит звездочка» и ««Дружные 
ребята»), дела пионерской организации, в которую они  
вступают в 3-м классе. Юные зрители легко отзываются на  
музыкально-развлекательные передачи, пишут в ««Будильник»», много 
почты получает веселая познавательная передача ««Турнир 
любознательных», хотя она и предназначена для более старших 
ребят. 
Учащиеся начальной школы смотрят многие пионерские 
передачи (««Пионерия», ««Костер»), о Родине, о событиях и  
приключениях детей — героев книг и передач. 
Одна из особенностей учащихся 1—2-х классов состоит в 
том, что у них формируются познавательные интересы в  
основном в пределах учебного предмета, представления об  
окружающей действительности более ориентированы, познание и 
деятельность для них приобретают большую значимость, если 
она признается полезной для других людей. Интересы не  
отличаются постоянством и по-прежнему в значительной степени 
связаны с уже имеющимся жизненным опытом. 
73 
Отдельные передачи направлены на то, чтобы укрепить 
любовь детей к чтению, а родители и педагоги должны всячески 
поощрять просмотр подобных передач («Кот-книгоноша», «По- 
лочка-читалочка» и другие). Передачи пытаются увлечь ребят 
коллективной деятельностью, подсказывают им совместные 
дела (уход за цветами в живом уголке, изготовление подарков- 
игрушек для детского сада и т. п.). 
Ребята в начальной школе многие моральные категории 
понимают еще неточно, смелого человека определяют как  
сильного, доброго как хорошего. Овладение нравственными  
понятиями не должно исчерпываться только понятием о них, без опоры 
на эмоциональное восприятие, переживание. Поэтому  
родители, бабушки и дедушки, сидящие вместе с детьми у телевизора, 
не должны быть нейтральными или назидательными. Они 
должны живо реагировать на события и поступки, которые 
показываются по ТВ, высказывая свое отношение к ним, а 
также поддерживать переживания своих детей. 
Следующая группа юных зрителей, пожалуй, самая  
многочисленная, самая активная и с готовностью отдающая  
свободное время своему другу- телевизору,— это подростки 10—13 
лет. 
Главная психологическая пружина, прочно прижимающая 
их к стулу, дивану, креслу, — это огромная познавательная 
потребность, которую можно охарактеризовать словами «Хочу 
все знать». Их буквально интересует все: война и муравей, дно 
океана и операция на сердце, луч лазера и техника  
стихосложения, героические деяния людей и пионерские повседневные 
заботы. Если дошкольников и младших школьников  
интересовало в первую очередь то, что связано с уже знакомыми 
вещами, понятиями, окружающими их в повседневной жизни, то 
подростков интересуют предметы, явления, о которых они не 
имеют ясного представления. 
Если у малышей сами вещи вызывают интерес, у младших 
школьников — занимательная форма, сюжетность передачи, то 
подростки сосредоточивают внимание на заинтересовавшем их 
явлении, событии. Они любят сравнивать новые знания с ранее 
74 
полученными. Сказка, фантазия по-прежнему еще привлекает 
их, однако уже больше нравится ««серьезная» передача  
(например, документальные фильмы «Клуба кинопутешествий»» и т. п.). 
В этом возрасте ребята любят коллекционировать,  
систематизировать, проникать в существо, докапываться до истины. У них 
нет еще устойчивых интересов. Они с удовольствием смотрят 
передачу по истории русского флота и о правилах уличного 
движения. 
Родителям надо не укорять ребят, что они ничем  
постоянным не интересуются, меняют свои увлечения. Наоборот,  
период поисков надо сделать насыщенным. Телевидение как раз и 
стремится широко распахнуть двери и окна в мир, как можно 
шире познакомить детей со всеми отраслями человеческой 
деятельности, с природой, пробудить в них жажду познавать, 
интерес к познанию. Только так можно действительно помочь. 
Узнав многое, легче найти то, что ближе сердцу и их  
способностям. 
Другая задача, которую ставит ТВ,— углублять и развивать 
интересы тех ребят, у которых они уже проявились, к чему-то 
определенному. 
Передач для подростков на всех студиях создается много. 
Циклы, серии, рубрики нередко меняют названия, но  
педагогическая сущность вещания для любой возрастной группы  
остается прежней. (Именно поэтому мы называли не только  
сегодняшние передачи, но и те, которые были раньше). Однако 
содержание передач для подростков ясно угадывается в  
названиях рубрик: «Пионерия», ««Край родной», ««Священные места 
нашей Родины», «Зеленая ракета», «Зарница», «Сами о себе», 
««Всемирный следопыт», «Турнир умелых», «Для  
любознательных нет наук незанимательных», «Школьникам об искусстве», 
«Читай- город», «Путешествие в книгу», «Спортивная юность», 
«Веселые старты», «Часовые природы», «Встречи юнкоров с 
интересными людьми». 
Если же по названию вы затрудняетесь представить  
содержание передачи, а аннотацию в печатной программе прочитать 
не смогли, то достаточно посмотреть одну-две передачи — и 
многое проясняется. 
75 
Как правило, передачи одной рубрики, как звенья одной 
цепи, связаны между собой общей проблематикой, целью и 
смыслом. В то же время каждая передача имеет  
самостоятельный законченный характер по форме и содержанию. Вспомните 
«Будильник», «Веселые старты», «Приходи, сказка»,  
«Искатели», «Олимпиады», популярную передачу телевидения ГДР — 
«Делай с нами, делай как мы, делай лучше нас». 
Создатели цикла заботятся об опознавательных знаках — 
постоянные графические заставки и музыкальные позывные и 
часто — постоянные ведущие. 
Все это очень помогает зрителю ориентироваться в  
программе и заранее наметить передачи для просмотра. 
Но вернемся к нашему разговору о типах передач по 
возрастам. Итак, передачи для старших школьников — ребят 
8—10-х классов. 
Передач для юношества создается меньше, чем для  
малышей и подростков. Объясняется это тем, что старшеклассникам 
доступны и интересны многие взрослые программы. Да и в 
школе они уже в 9-м классе знакомятся с такими сложными 
многоплановыми произведениями, как «Война и мир», «Анна 
Каренина» Л. Н. Толстого, «Гроза» А. Н. Островского, романы 
Достоевского; в спецшколах им рекомендуют читать в  
подлиннике произведения Шекспира и Б. Шоу. 
Тематика взрослых телепередач становится и их тематикой. 
Разумеется, и в этом возрасте еще не снята проблема  
ограничения просмотра не столько из-за содержания передач, 
сколько из-за необходимости соблюдения старшеклассниками 
режима отдыха. Хотя совершенно игнорировать вопрос о 
содержании передач тоже не следует. Не следует потому, что в 
поздний вечерний час показываются фильмы и спектакли,  
которые в большинстве случаев посвящены проблемам семейной 
жизни или другим вопросам, к которым вряд ли стоит раньше 
времени привлекать внимание подростков и тем самым  
искусственно возбуждать интерес к этим темам. 
Юношество. Переплетение взрослости и детства. Это 
период, когда люди пытаются понять себя, их волнует вопрос: 
76 
«Какой я?», т. е. внутренний мир человека, становление  
характера, сознательное и целенаправленное самовоспитание, 
выработка жизненных идеалов, взглядов на жизнь, проблемы 
любви, взаимоотношений между людьми; определяется и  
профессиональная направленность, попытки к самоутверждению, 
повышается интерес к поэзии, политике, психологическому 
роману. Они с любопытством смотрят молодежные программы, 
стремятся вникнуть в проблемы, которые их ждутзавтра,  
пытаясь определить к ним свое отношение и представить свое  
будущее. 
Самой популярной передачей среди юношества, как  
показали опросы, был КВН — клуб веселых и находчивых.  
Очевидно, он больше всего соответствовал психологическим  
особенностям старшеклассников, создавал у них оптимистический 
взгляд на свое будущее, ибо КВН привлекал не только своим 
обаянием, веселостью, остроумием, но и тем, что незаметно 
создавал собирательный портрет современного молодого,  
пытливого, граждански-активного человека с широким  
диапазоном интересов, коллективиста, честного, прямого и  
доброжелательного, волевого и целеустремленного. Эта передача  
давно имеет филиалы на студиях, а также 8 институтах,  
молодежных коллективах предприятий, совхозов. 
В театре «Публицист» старшеклассники как бы все вместе 
встречаются со сложной современной жизнью, с нелегкими 
судьбами, сами судят, сами отвечают на вопросы, формируют 
свое отношение к жизни. 
Но ни молодежные, ни взрослые программы не могут  
исчерпать интересы, свойственные юношескому возрасту. 
Молодые рабочие, учащиеся техникумов и ПТУ, студенты — 
они уже нашли главное в жизни — свой трудовой путь, свою 
профессию, часто и призвание, а старшеклассникам это еще 
предстоит. 
Не случайно в циклах и передачах для юношества («Тебе 
юность», «Рассказы о профессиях», «Мое место в жизни»), в 
отдельных тематических передачах о школьных производствен- 
77 
ных бригадах и т. д. в центре внимания проблема поиска своего 
жизненного пути. 
Можно выбрать профессию, но это далеко не все. Юношам 
и девушкам предстоит ответить не только на вопрос «кем 
быть?»», но и «каким быть?". 
Нередко телестудии создают передачи ««Сделать жизнь 
с кого?», названные строчками из известной поэмы В.  
Маяковского. Задав такой вопрос, Маяковский отвечал: «Делай 
ее с товарища Дзержинского». Кристально чистая жизнь  
Феликса Эдмундовича, сгоревшего в борьбе за идеалы революции, 
прекрасный пример для юношества. Маша революционная 
история богата героями, которые отдавали и отдают свою жизнь 
без остатка за счастье народа. 
Различные студии телевидения создали в своих передачах 
замечательную галерею портретов людей, чьи имена  
составляют боевую и трудовую славу нашей Родины (циклы ««Жизнь 
замечательных людей», «Они были первыми» и др.). 
Старшеклассники, проявляя огромный интерес к  
нравственным проблемам, отбирают из телеспектаклей и фильмов 
те моральные ««кирпичики», из которых они строят свой  
характер. Утверждение нравственных ценностей, защита их  
способствует становлению общественной активности, проявлению в 
практической жизни непримиримого отношения к недостаткам, 
готовности выполнить свой гражданский долг перед Родиной. 
Телестудии систематически готовят передачи, освещающие 
опыт работы школьных комсомольских организаций, ставящих 
проблемы взаимоотношения коллектива и личности. 
Следует особо остановиться на передачах для юношества, 
которые идут по Центральному телевидению более десяти лет 
— это цикловые передачи — ««Клуб «Искатели» и ««Олимпиады». 
Клуб ««Искатели» во многом способствовал развитию  
«искательского» движения. Он обратился к туристам с просьбой свои 
походы посвящать сбору материалов для школьных и  
краеведческих музеев. По всей стране были созданы клубы и кружки 
искателей, которые добровольно брали на себя темы,  
предлагаемые телевизионными студиями, и организовывали поиски 
78 
документов, материалов, экспонатов для школьных и  
краеведческих музеев. Кроме студий телевидения, к ребятам с  
заданиями обращались комсомольские организации, советские и 
научные учреждения, колхозы и совхозы, писатели и  
журналисты. Это были различные задания, начиная от изучения 
режима малых рек и поиска местного строительного  
материала до сбора документов и экспонатов, отражающих 
боевую и трудовую славу советского народа. 
Школьные олимпиады, которые проводятся по многим 
отраслям знаний — математические, биологические,  
физические, химические и др., собирают у экранов телевизора  
большое количество пытливых юношей и девушек. В одиночку и 
целыми классами (а иногда и школами) ищут они ответы на 
вопросы, поставленные в передачах, делают расчеты, чертежи 
моделей. Письма ребят на студию полны творческой энергии и 
фантазии. Например, какой смелостью (а часто и зрелостью) 
конструкторской мысли радовали работников телевидения и 
членов жюри ответы на вопрос: «Сделайте и защитите свой 
проект и чертеж будущего самолета Арктики». Подготовка  
ответов 'потребовала от ребят не только вдумчивой работы с 
научной литературой, но и консультаций со специалистами, 
посещения музеев, соответствующих учреждений. 
Самое отрадное, что в олимпиадах принимают участие 
школьники из далеких сел и поселков. Олимпиады не только 
углубляют интерес у ребят к тому или иному предмету, но и 
требуют самостоятельной и активной работы. Непременные 
участники этих передач — члены жюри и консультанты, — 
крупные ученые. Разбором писем занимаются аспиранты и  
студенты старших курсов. Школьники вводятся в мир  
действительно актуальных научных проблем и практического их  
претворения. Десятки победителей третьего тура вызываются в 
Москву, где ученые-практики, покорители Севера или морских 
глубин, инженеры и летчики, рабочие-рационализаторы 
вручают дипломы победителям Олимпиады. 
Олимпиады также помогают выявлять талантливых и  
способных ребят для пополнения вузов, а путевкой, помогающей 
79 
им поступить в институты, служит диплом участника. Но самое 
главное, участвуя в олимпиаде или подобных передачах (в Горь- 
ковской студии телевидения — Краеведческие викторины, в 
Риге, например, — «Твой проект автомобиля 2000 года» и т. п.), 
ребята проходят школу самостоятельного пытливого труда. 
Мы намеренно так много строк посвятили этим двум  
передачам для юношества. Нам хотелось обратить на них особое 
внимание родителей, чтобы они советовали своим  
взрослеющим детям включаться в подобные передачи, ибо они  
побуждают юношей и девушек к активной и полезной деятельности, 
помогают им ««найти себя», приучают их самостоятельно  
овладевать знаниями. 
Студии некоторых городов (например, Ленинграда) готовят 
для старшеклассников специальные турниры между школами и 
в том числе по проблемам литературы и искусства, музыки. И 
хотя утихли споры о физиках и лириках, и доказано, что полет 
мысли, свойственный поэтам, еще более оттачивает острый 
точный ум физиков и обогащает их фантазию, нет-нет да и 
встречаются юноши, которые с максималистской крайностью 
утверждают, что для жизни поэзия и искусство не нужны.  
Поэтому, дорогие родители, всячески поддерживайте интерес  
ваших детей к литературным передачам для молодежи. 
Мир красоты и поэзии должен стать миром ваших детей. 
Телевидение создает передачи для молодежи, которая уже 
выбрала свою профессию — для учащихся ПТУ и техникумов. У 
этой группы ребят имеются свои специфические интересы. 
Передачи для них называются «От 14 до 18». Они своей  
серьезностью в постановке производственных вопросов, проблемами 
овладения профессиями и подготовкой к трудовой жизни во 
многом привлекают и учащихся 8-10-х классов. Своим  
отношением к будущему, к труду, к жизни, деловитостью,  
целеустремленностью учащиеся профтехучилищ заражают  
старшеклассников. 
Поэтому очень полезно школьникам 8—10-х классов  
смотреть передачи, предназначенные для учащихся профессио- 
80 
нально-технических училищ, а также для молодых рабочих,  
колхозников и других специалистов народного хозяйства. 
Чем закончить разговор о передачах, о программах? 
Пожалуй, передачей, о которой мы пока еще ничего не 
сказали и которую ежедневно смотрят дети 4, 5, 6, 7, 8 и 9 лет 
по второй программе Центрального телевидения — это  
«Спокойной ночи, малыши». 
Бабушки и дедушки, вы, наверно, облегченно вздохнули: 
теперь вам не надо выдумывать перед сном сказку, теперь 
малыши слушают очередную вечернюю сказку по ТВ. Но  
малышам оказалось этого мало, они просят вас комментировать 
вечерние сказки, а иногда и писать на студию письма: «Сказка 
очень понравилась, видели уже два раза, повторите еще». 
Смысл вечерних детских передач ясен: выработать у детей 
устойчивую привычку ложиться спать в определенное время. 
Сказка и до телевидения была нечто вроде точки в конце  
предложения, своего рода обозначением, что день окончен и пора 
спать. 
Хочется, чтобы взрослые готовили ребят к этой передаче. 
Пусть дети сначала почистят зубы и умоются на ночь, а потом 
уже усядутся у телевизора. Сказка кончена — и сразу же 
малыш должен ложиться в постель. Не «разгуливайте»» его 
новыми разговорами, не будите его фантазию, пусть он под 
впечатлением виденного засыпает. 
Работники телевидения стремятся к тому, чтобы сказки не 
были страшными, будоражащими ребят, чтобы они  
заканчивались спокойным «счастливым» концом. 
Итак, программы телевидения — это не случайный набор 
передач, это части единого целого. 
Просмотр отдельно взятой передачи возможно не всегда 
оказывает заметное воспитательное воздействие. Но  
целенаправленный и сознательный подход к выбору передач  
постепенно накапливает у юного зрителя определенные  
представления, вырабатывает суждения и, таким образом, способствует 
формированию взглядов. 
Зная своих детей, родители, несмотря на все рекомендации, 
81 
сами должны разумно регулировать просмотры телепередач. 
Иногда, в целях создания мирного домашнего климата, в 
целях укрепления внутрисемейных отношений, вполне  
возможно, что родители вместе с детьми будут не часто, но регулярно 
смотреть взрослые передачи. 
Это могут быть музыкальные, географические,  
театральные, спортивные программы. Детей 9—10 лет вряд ли стоит 
приобщать к передачам продолжительностью более 30 минут. 
Они устанут и будут просто «вертеться» между взрослыми, 
отвлекать их. А детям 12—13—14 лет уже можно иногда, один 
раз в неделю, вместе со взрослыми посмотреть (до 9 часов) 
фильм или послушать концерт, просмотреть спортивную  
программу. 
Но это, так сказать, формальное признание того, что дети 
могут смотреть передачи для взрослых. 
А если говорить по существу, то мы даже отстаиваем право, 
полезность и необходимость просмотра взрослых программ 
детьми. 
Мы исходим из того, что правильное воспитание происходит 
тогда, когда дети не изолируются от старших, когда каждый 
ребенок имеет возможность усваивать знания, опыт, нормы 
поведения от взрослых и, в свою очередь, передавать свой 
опыт и знания младшим. Можно сказать, что социальное  
развитие личности ребенка происходит в результате ориентации на 
взрослых, усвоения их ценностей и норм поведения. 
Поэтому неправомерна постановка вопроса: хорошо или 
плохо, что дети смотрят взрослые передачи? Однозначного 
ответа быть не может. Во-первых, о каких детях идет речь — 
четырехлетних или четырнадцатилетних? Во-вторых, о каких 
передачах идет речь? Если о фильме «Чапаев» или фигурном 
катании на коньках, то мы можем с уверенностью  
рекомендовать их для юных зрителей. А такой фильм, как «Перевод с 
английского»», вызвал добрую реакцию пятиклассников и  
десятиклассников, студентов и учителей. И каждый зритель нашел 
что-то ценное для себя. 
82 
Конечно,есть «темы детские» и детские передачи, есть 
«взрослые темы» и передачи для взрослых. 
Но гораздо больше в жизни таких ситуаций, которые  
охватывают широкий круг вопросов моральных,  
общественно-политических, эстетических, которые интересуют и взрослых, и 
детей. Каждый видит в них то, что больше отвечает его  
потребностям и пониманию. 
В этом случае скорее встает вопрос об ограничении  
просмотров вообще, чем о категорическом запрещении передач 
до... 16 лет. Мы за контроль, за разумное регулирование 
просмотров, и за запрещение только тех передач, в которых 
делается акцент на интимные отношения между мужчинами и 
женщинами. Где вообще ставятся сугубо взрослые проблемы. 
Причем, такие передачи в основном идут после 21 часа 30 
минут, когда дети младшего и среднего школьного возраста 
должны укладываться слать. 
К сожалению, в практике работы телевидения еще бывает, 
когда художественные программы «только для взрослых» 
включаются в ранние часы. 
Например, за неделю с 8 января по 14 января 1973 года до 
21 часа прошли три программы, которые были отнюдь не  
предназначены для просмотра детьми младшего и среднего  
школьного возраста. В частности, во французском фильме «Гром 
небесный» рассказывается, как герой фильма Брасак добрый и 
сильный человек, но в то же время пьяница и буян, привел в дом 
проститутку и «помог ей встать на ноги». Его речи вовсе не 
предназначены для ушей и глаз подростков. Порой он груб, 
прямолинеен и циничен. Многие родители этот фильм не видели 
раньше и потому разрешили его смотреть детям. А потом  
испытали неприятные минуты, когда по ходу фильма приходилось 
либо выключать телевизор, либо отсылать ребят в другую  
комнату. 
В то же время на этой неделе после 21 часа 30 минут  
показывались такие фильмы, как «Котовский», комедия «Укротители 
велосипедов», трилогия по роману А .Н. Толстого «Хождение 
83 
по мукам», ряд хороших музыкальных программ и  
документальных фильмов. 
Знакомство с программой поможет избежать во многих  
случаях подобных ошибок. 
««Взрослые» передачи, такие, как «Клуб кинопутешествий»», 
«Клуб «Подвиг» и «Клуб «Поиск», телевизионные новости,  
многие спортивные программы, музыкальные, классические  
литературные произведения, выступления политических  
комментаторов, советские кинокомедии, нисколько не противопоказаны 
детям. Но при этом надо в каждом отдельном случае  
руководствоваться разумом, т. е. помнить, что у каждого возраста есть 
«предел вместимости информации», за которым приходит  
утомление и эмоциональное снижение восприятия. 
А теперь несколько слов о развлечении. Не знаю почему, но 
у некоторых людей установилось предосудительное отношение 
к развлечению, как к чему- то недостойному серьезного  
человека, а в отношении детей — как к баловству, бессмысленной 
забаве, пустой трате времени. Иной взрослый человек,  
просиживающий часами у телевизора, ни за что не скажет «я  
развлекаюсь», ему как-то неловко развлекаться, и он подберет другие 
слова — «я отдыхаю». 
Однако если мы вернемся к анализу передач, которые при 
исследовании вошли в рубрику, «использую, ТВ как средство 
отдыха», то получится, что и взрослые, и дети отдыхают, в 
первую очередь просматривая мультфильмы, комедии  
(кинофильмы и спектакли), «Кабачок 13 стульев», «Артлото», «Алло, 
мы ищем таланты», «Будильник», «А ну- ка, парни», «Турнир 
любознательных», фигурное катание, хоккей, а также  
вспоминают передачи, которые были раньше — «КВН», «Голубой  
огонек» (по субботам), «Терем-теремок», «С днем рождения», 
«Сердитые странички» и др. А ведь основная цель этих передач 
— развлекательная. 
Не надо развлечение уличать в «неблаговидном  
поведении», обзывая «развлекательством». Развлечение — это 
весьма эффективный способ снятия нервного напряжения и 
восстановления сил, затраченных в течение дня. 
84 
И если девочка Таня пишет: «Я люблю смешные истории и 
песенки», если Слава сообщает, что он заливается смехом, 
когда смотрит мультфильмы, хотя и не маленький —  
шестиклассник, а девятиклассник Толя просит создавать веселые 
комедии и телеконцерты, так как они поднимают настроение и 
помогают с юмором и оптимизмом относиться к разным  
неприятностям, то и телевидение и родители должны к этому  
прислушаться. Итак, развлечение —дело необходимое, серьезное. 
И в заключение —два совета. 
Детскому телевидению — быть «повеселее»». Создавать для 
детей передачи, наполненные приключениями, музыкой,  
шуткой, игрой, не уступающие некоторым взрослым передачам. 
Родителям — радуйтесь, когда ваши дети заливаются  
смехом у экрана, когда у них блестят глаза от удовольствия, но 
помните, что получение развлечений не должно стать у ваших 
детей самоцелью, устойчивым стремлением, привычкой. 
Человек, привыкающий к неумеренному потреблению  
удовольствий, может сформироваться неполноценной личностью, 
у него происходит смещение понятий о подлинных ценностях — 
о смысле труда, жизни, представление о счастье. 
ВМЕСТО ЗАКЛЮЧЕНИЯ 
Мамы и папы, бабушки и дедушки! 
Вы знаете своих детей, их увлечения, привязанности,  
предпочтения лучше, чем кто-либо. И все же мы хотим предложить 
некоторые приемы как бы «самопомощи» в выявлении роли ТВ 
в жизни ваших детей, в развитии их интересов. 
Первый такой прием — это ведение учащимися  
специальных дневников, в которые они в течение месяца ежедневно 
записывают все, что они делают после прихода из школы. Это 
своего рода самофотография дня. 
Форма дневника проста, но родителям придется регулярно 
следить за его заполнением. 
Кроме этого приема-анализа, применяют еще и другие. 
Например, ребятам 5—7-х классов говорили: 
«Если бы вас сделали на одну неделю временно  
исполняющим обязанности директора телевидения, то какую бы Вы 
составили программу с двух часов дня и до одиннадцати  
вечера?» 
Ребята охотно брались за эту работу, представляя себя и 
взрослым, и директором. В программах, которые они  
составили, ясно выразились их предпочтения. Во-первых, как правило, 
на 70 процентов это были — кинофильмы и передачи  
развлекательного характера, и только на 10 процентов познавательного 
и информационного. Однако каждая программа отдельно 
давала прекрасный материал о месте и роли ТВ в жизни 
ученика, об отношении к книге, спорту и т. д. «Детские 
программы»» давали ключ к беседам, а главное к комплексу 
86 
других воспитательных совместных действий семьи и школы 
для пробуждения у ребенка интересов, воспитания культуры 
избирательного отношения к средствам массовой информации, 
а также организации творческой деятельности на основе  
передач. 
Мне кажется, что может принести пользу еще один опыт, 
который используется время от времени в одной сельской и 
московской школах. Ребятам дают газету с опубликованной 
недельной программой, но после того, когда неделя уже 
прошла. Предлагается отметить кружочком те передачи,  
которые понравились, квадратиком — которые можно было бы и не 
смотреть (т. е. передачи, которые не принесли  
удовлетворения) и треугольником —то, что активно не понравилось. 
Обычно анализ такой своеобразной анкеты дает  
возможность выяснить, какого числа и сколько было просмотрено  
программ, во-вторых, сколько времени потрачено на просмотр 
передач, которые ничего не дали, говорят о зря потерянном 
времени. Передачи, которые отмечены как хорошие или плохие, 
дают основание судить о вкусах и направленности интересов. 
Последующие наблюдения и беседы с учениками позволяют 
родителям и учителям более осмысленно направлять  
зрительскую деятельность детей. 
В заключение хотелось бы приложить анкету, которую вы 
можете провести в своей семье и которая наверняка наведет 
вас на педагогические размышления. 
Анкета 
(Нужный ответ «да», «нет» обвести кружочком справа) 
1. Как ты чаще всего проводишь свободное время? 
Гуляю, играю, читаю, смотрю телевизор, слушаю радио, 
занимаюсь любимым делом (каким: спорт, техника и т. п.),  
занимаюсь пионерской (комсомольской) работой. 
87 
2. Из того, что содержится 8 первом вопросе, чему бы ты 
хотел отдавать больше времени? 
3. А чему ты отдаешь предпочтение на самом деле? . . . 
4. Выбираешь ли ты заранее телепередачи? 
5. Кто тебе чаще помогает выбирать передачи? 
(Отец, мать, дедушка, бабушка, учитель, товарищ, печатная 
программа, никто не помогает). 
6. Постарайтесь припомнить, помогли ли какие-нибудь  
телепередачи подготовиться к заданному уроку 
заинтересовать какой-либо профессией 
выбрать интересную книгу 
изменить'мнение, взгляды, намерения 
убедить кого-то в своей правоте 
найти интересное дело 
увлечься общественными делами 
7. Служат ли просмотренные передачи темой для  
разговоров: 
со сверстниками 
со взрослыми 
8. Какие передачи тебе больше всего нравятся (кроме  
фильмов и спектаклей)? 
9. Твоя любимая детская передача (назови) 
10. Твоя любимая взрослая передача (назови) 
Почему? 
88 
11. Мешает ли тебе телевидение учить уроки? 
часто мешает 
иногда мешает 
не мешает 
12. Считаешь лы ты, что твой режим дня правильно  
учитывает время для уроков, прогулок и для просмотра телепередач? 
Приложение 
ИЗ ИСТОРИИ СОВЕТСКОГО ТЕЛЕВИДЕНИЯ 
Родоначальник телевидения русский ученый 
В. Л. Розинг вряд ли представлял себе реально, как будет  
выглядеть его изобретение через несколько десятков лет. 
Идея изображения на расстоянии с помощью  
электромагнитных волн возникла у него в 1907 году, а 22 мая 1911 года он 
впервые в мире демонстрировал на экране своей  
электроннолучевой трубки темные линии на светлом фоне. 
Однако' серьезная работа начала проводиться после 
Октябрьской революции в Нижегородской радиолаборатории, 
созданной по инициативе В. И. Ленина и руководимой ученым 
революционером М. А. Бонч-Бруевичем. 
Уже в 1921 году В. И. Ленин, просматривая очередное  
сообщение о положении дел в фотолаборатории, обратил внимание 
на одну фразу, в которой говорилось, что «получен новый  
фотоэлемент» и «что при соответствующем усовершенствовании он 
позволит видеть на экране подвижное изображение человека 
при радиотелефоне...»» 
Владимир Ильич оценил важность и перспективность этого 
изобретения и распорядился, чтобы лаборатории была оказана 
помощь в усовершенствовании этого прибора. 
90 
Чегез десять лет 1 октября 1931 года в Москве на ул. 25 
октября в доме № 7 началось телевизионное вещание.  
Изображение принималось в Одессе, Харькове, Нижнем Новгороде, 
Томске и других городах Советского Союза. Экраном служила 
электроннолучевая трубка, созданная советским ученым С. И. 
Катаевым. 19 октября телевидение демонстрировалось в 
Ленинграде. 
Однако до 15 августа 1932 года изображение по-прежнему 
оставалось неподвижным и немым. Конечно, передатчики были 
несовершенными, а экран был чуть больше спичечной коробки. 
25 марта 1938 года Московская студия телевидения впервые 
показала кинофильм «Великий гражданин»». 
Война помешала нормальному развитию телевидения, 
однако советское телевидение первое в Европе возобновило, 
свои передачи в 1945 году. 
16 июня 1949 года вступил в эксплуатацию новый  
телевизионный центр на ул. Шаболовка. 29 июня Московское  
телевидение провело первую трансляцию футбольного матча со  
стадиона «Динамо»». 
12 апреля 1961 года в космосе советский гражданин Юрий 
Гагарин. В течение ряда дней советское телевидение живет 
этим событием; восторженная встреча москвичей Юрия  
Гагарина, прием в Кремле, Международная пресс-конференция, 
вечер на студии телевидения, выступления ученых и инженеров 
— участников космических исследований. Миллионы людей 
являлись свидетелями памятных событий. 
18 марта 1965 года — советские зрители на своих экранах 
увидели необычное зрелище — выход в открытый космос  
Алексея Леонова из космического корабля, пилотируемого Павлом 
Беляевым. 
Ноябрь 1967 года — вступил в строй гигантский  
телевизионный комплекс имени 50-летия Октября в Останкино (2600 
помещений, 170 тысяч кв. м., 21 павильон, из которых два — по 
1 тыс. кв. м.). 
Развитие телетехники продолжается. Изучается проблема 
передачи изображения с помощью лазерного луча, создаются 
звуковые установки со стереофоническим звучанием,  
кассетные передатчики типа миниатюрных видеомагнитофонов. 
В настоящее время телефикацией охвачена территория, на 
которой проживает более 70% населения Советского Союза. 




	
	
	\newpage
	\tableofcontents
	
	\thispagestyle{empty} % 
	
	\newpage
	
	\setcounter{secnumdepth}{0}
	
	\phantomsection
	
		\section*{Описание}
	
	{\bf Название:} А если в семье телевизор... и дети?
	
{\bf Автор:} Ксенофонтов Валентин Васильевич

{\bf Редактор:} О. Г. Свердлова
	
{\bf Издательство:} М.: <<Знание>>. 1973. - 96 с.
	
		{\bf Аннотация:} Ученые всех стран мира сходятся в одном — обилие телевизионной информации, даже полезной по своему содержанию, без правильного руководства со стороны взрослых воспитывает у детей душевную и нравственную пассивность, приучает довольствоваться малым, не говоря уже о том, что оно отвлекает от чтения, спорта и других занятий. И в то же время телевидение представляет детям весь мир наглядно, в ярких динамичных образных формах. Чего только не увидишь по телевизору! Луноход и дно океана, сказку и экзотических животных. 
		
		Книга учит родителей правильно использовать огромные возможности телевидения для воспитания детей. Она рассчитана на самые широкие круги читателей. 
		


		\thispagestyle{empty} % выключаем отображение номера для этой страницы

	
\end{document}


